\documentclass[11pt,a4paper]{article}
\usepackage[T1]{fontenc}
\usepackage[utf8]{inputenc}
\usepackage{lmodern}
\usepackage[margin=26mm]{geometry}
\usepackage{amsmath,amssymb,amsthm}
\usepackage{booktabs,longtable,array}
\newcolumntype{L}[1]{>{\raggedright\arraybackslash}p{#1}}
\usepackage{microtype}
\usepackage{cite}
\usepackage[hidelinks]{hyperref}
\usepackage{url}
\urlstyle{same}
\setlength{\emergencystretch}{2em}
\setlength{\parskip}{3pt}
\setlength{\parindent}{1.2em}
\newtheorem{proposition}{Proposition}[section]
\newtheorem{corollary}[proposition]{Corollary}
\theoremstyle{definition}
\newtheorem{definition}[proposition]{Definition}
\newtheorem{assumption}[proposition]{Assumption}
\theoremstyle{remark}
\newtheorem{remark}[proposition]{Remark}
\newcommand{\TCGS}{TCGS--SEQUENTION}
\newcommand{\dd}{\mathrm{d}}
\newcommand{\ii}{\mathrm{i}}
\newcommand{\Tr}{\operatorname{Tr}}
\newcommand{\Var}{\operatorname{Var}}
\newcommand{\Cov}{\operatorname{Cov}}
\newcommand{\diag}{\operatorname{diag}}
\newcommand{\R}{\mathbb{R}}
\newcommand{\musp}{\mu'_s}
\newcommand{\nb}{\bar n}
\newcommand{\lop}{\ell_z}
\newcommand{\topc}{\tau_{\mathrm{op}}}
\newcommand{\doi}[1]{\href{https://doi.org/#1}{doi:\nolinkurl{#1}}}
\newcommand{\arxiv}[1]{\href{https://arxiv.org/abs/#1}{arXiv:\nolinkurl{#1}}}
\hypersetup{pdftitle={Vacuum-Fluctuation Imaging and Three-Dimensional Carrier Irreducibility},pdfauthor={Henry Arellano-Pena},pdfsubject={Effective dimensional reduction under TCGS-SEQUENTION}}
\title{Vacuum-Fluctuation Imaging and\\
Three-Dimensional Carrier Irreducibility\\[0.6em]
\large A Quantitative Empirical Anchor for the\\
Carrier--Representation Distinction under TCGS--SEQUENTION}
\author{Henry Arellano-Pe\~na\\\small Nuevo Estandar Biotropical NEBIOT S.A.S., Bogot\'a, Colombia\\\small\href{mailto:harellano@unal.edu.co}{harellano@unal.edu.co}}
\date{27 September 2026}
\begin{document}
\maketitle
\begin{abstract}
The imaging of ground-state spin fluctuations in a planar, coherently coupled $^{39}\mathrm K$ Bose--Einstein condensate provides a precise setting in which physical carrier dimension, effective field dimension, measurement dimension, and source ontology can be separated. This paper interprets that experiment under the governing \TCGS\ architecture of non-temporal four-dimensional Counterspace, a three-dimensional physical shadow, and selector-indexed 4D-Counterduction. Its central empirical anchor is the realization of effective two-dimensional quantum-field behaviour by a finite-width atomic carrier, rather than an experimental determination of universal dimensional ontology. We derive the transverse-mode reduction, the confinement dependence of the effective interaction, and a finite-energy obstruction to exactly zero transverse width within the massive-atom quantum model. The reported transverse trap frequency of approximately $1\,\mathrm{kHz}$ implies a nominal oscillator length $\lop\simeq0.509\,\mu\mathrm m$ and root-mean-square width $\sigma_z\simeq0.360\,\mu\mathrm m$. The suppressed direction remains physically registered in the coupling $g_{2D}\propto\lop^{-1}$ and in the transverse form factor. Independently, a canonical covariance analysis recovers the measured vacuum-dominated $\omega_k^{-1}$ fluctuation scaling, distinguishes it from thermal scaling, and shows why the measured covariance anisotropy is not a complete covariance invariant or a dimensional estimator. We audit detection efficiency, mode selection, adiabaticity, and the separation between stationary ground-state uncertainty and deliberately driven covariance rotation. A non-identifiability construction demonstrates that equal planar observables need not identify the transverse carrier, while an axial-response protocol can restore the missing information. The resulting framework interpretation preserves effective two-dimensional physics without promoting its reduced coordinates into an autonomous lower-dimensional physical substrate. Comparisons with spectral dimensional reduction and holographic duality retain the technical meaning of each construction. The paper closes with transferable confinement, covariance-tomography, and held-out calibration tests that make the carrier--representation distinction operationally explicit.
\end{abstract}
\noindent\textbf{Keywords:} \TCGS; three-dimensional carrier irreducibility; vacuum fluctuations; quasi-two-dimensional Bose gas; transverse confinement; effective field theory; covariance tomography; source--shadow geometry; 4D-Counterspace; 4D-Counterduction; operational time; dimensional reduction; invariants; artifacts.

\section{Introduction}
\label{sec:intro}
An image of a fluctuating quantum field invites an immediate question: what, exactly, has been imaged? Zhang and collaborators answer this experimentally by encoding a bosonic field in the spin degrees of freedom of a coherently coupled, two-component $^{39}\mathrm K$ Bose--Einstein condensate (BEC). Their images resolve spatially structured measurement outcomes whose scale-dependent amplitudes agree with the ground-state covariance of the effective spin field. The experiment is therefore a measurement of a prepared many-body system, not a photograph of metaphysical nothingness. Its vacuum is defined relative to the Hamiltonian and mode decomposition of that system.\cite{zhang2026}

The distinction matters especially when the field is called two-dimensional. The effective variables depend on two in-plane coordinates, but the atoms are confined in the third direction by a finite harmonic potential. The reduced description does not remove the transverse wavefunction: it suppresses or integrates out selected transverse degrees of freedom. The experimentally specified confinement supplies a quantitative route for distinguishing the dimensionality of the effective description from that of the physical carrier. This distinction is already intrinsic to the theory of quasi-two-dimensional atomic gases.\cite{zhang2026,hadzibabic,petrov}

Under \TCGS, the observable physical domain is a three-dimensional shadow carrier, whereas complete source content belongs to non-temporal four-dimensional Counterspace. The present application begins from that governing architecture. It asks how a particular laboratory reduction should be typed within it, rather than attempting to reconstruct the source ontology from a single fluctuation spectrum. The framework's distinction between a source configuration, a constitutive relation, an operational carrier, and a mathematical representation is essential to the analysis.\cite{foundation,logical}

The central thesis has a restricted but substantive content: \emph{a successful two-dimensional quantum-field description can retain measurable dependence on a physically three-dimensional carrier}. Here the dependence is explicit in the transverse oscillator length and the effective interaction. This makes the experiment an empirical anchor for the carrier--representation distinction. The broader claim that every admissible shadow carrier is three-dimensional remains part of the governing framework; its scope must not be silently replaced by a theorem about one atomic confinement Hamiltonian.\cite{foundation,planck}

This paper adds a quantitative synthesis rather than a competing fit to the experimental data. It connects the reported confinement to a nonzero transverse width, derives the consequences of reducing the atomic field to a planar sector, and specifies the finite-resource meaning of carrier irreducibility. It then rederives the spin-mode covariance and identifies which properties survive canonical readout changes. Finally, it distinguishes the experiment's genuine quantum result from three stronger claims that its observables do not select: universal source-side randomness, a literal zero-thickness physical substrate, and an ontological interpretation of the laboratory time coordinate. The empirical starting point throughout is the authors' reported apparatus, equations, protocols, and uncertainties.\cite{zhang2026}

The ordering of the argument is important. Sections~\ref{sec:evidence}--\ref{sec:types} establish the evidence and framework vocabulary. Sections~\ref{sec:reduction}--\ref{sec:numerics} derive the carrier result. Sections~\ref{sec:spin}--\ref{sec:metrology} analyse the covariance and its measurement. Sections~\ref{sec:time}--\ref{sec:comparisons} delimit the source-level interpretation and compare other dimensional constructions. Section~\ref{sec:tests} supplies a test programme. The appendices provide derivational details that can be checked without adopting the framework ontology.

\section{Experimental evidence and source-context}
\label{sec:evidence}
\subsection{What the atomic experiment represents}
The condensate contains two coherently coupled hyperfine components, denoted $\lvert\uparrow\rangle$ and $\lvert\downarrow\rangle$. In the convention of the experimental paper,
\begin{equation}
 \psi_{\uparrow,\downarrow}
 =\sqrt{\frac{n(1\pm Z)}{2}}
   \exp\!\left[-\ii\left(\theta\pm\frac{\phi}{2}\right)\right].
 \label{eq:spinor}
\end{equation}
The total areal density is $n$, the common phase is $\theta$, the population imbalance is $Z$, and the relative phase is $\phi$. The spin field is the effective field organized by the latter two variables. It is not the three-dimensional microscopic atomic field before transverse reduction, and it is not the electromagnetic vacuum.\cite{zhang2026}

The in-plane trap is approximately square, with side $L\simeq33\,\mu\mathrm m$. Along the transverse direction, the reported confinement is harmonic, with $\omega_z/(2\pi)\simeq1\,\mathrm{kHz}$. The mean spin- and density-interaction scales are approximately $\mu_s/h\simeq900\,\mathrm{Hz}$ and $\mu_d/h\simeq120\,\mathrm{Hz}$, respectively. These are rounded apparatus parameters, not exact constants. The numerical analysis below retains that precision status.\cite{zhang2026}

Homogeneous box traps and transverse confinement are established experimental constructions, rather than inventions of the present interpretation. The optical-box literature documents the control of bulk uniformity and boundaries, while experiments on transverse condensation explicitly distinguish confined motion from the emergence of in-plane coherence. These sources provide the apparatus context for the present carrier audit.\cite{navon,chomaz}

The immediate precursor to the vacuum-imaging study demonstrates the same platform's ability to encode massive relativistic spin fields and non-perturbative relative-phase structures. Its sine-Gordon correspondence concerns the long-wavelength sector of the relative phase, not a claim that the atomic gas is fundamental relativistic spacetime. The relative-phase domain-wall construction has a separate theoretical history in two-component condensates.\cite{precursor,son}

\subsection{Amplification and direct-readout protocols are different}
In the first quantitative protocol, the spin sector is prepared deep in the Rabi regime and the coherent coupling is suddenly reduced. The quench changes the quadratic Hamiltonian. Initial ground-state uncertainty is thereby converted into an enhanced, periodically varying phase signal. Figure~2 of Zhang et al. combines approximately 70 independent images per delay to extract the fluctuation structure factor. The resulting oscillation amplitude has a known dependence on the final coupling and wavevector.\cite{zhang2026}

In the second protocol, the coupling is reduced over approximately $9\,\mathrm{ms}$ to prepare the spin modes near their ground states at the selected final coupling. A subsequent rapid, strong-coupling rotation maps different quadratures into the measured spin projection. Figure~3 uses approximately 120 independent images per delay and obtains a fluctuation amplitude proportional to $\omega_k^{-1}$ over the retained range. A fit allowing thermal occupation gives $T=(6\pm6)\,\mathrm{nK}$. These are source-reported results, not estimates reconstructed here from the plotted points.\cite{zhang2026}

The distinction prevents two misreadings. First, the images at successive delays are not nondestructive frames of one persistent realization. Second, the periodic structure-factor signal is deliberately generated by a quench or readout rotation. The observation of ground-state uncertainty and the observation of a driven covariance rotation are related experimental achievements, but they are not the same physical statement.\cite{zhang2026}

\subsection{The meaning of ``direct'' and of ``vacuum''}
Here ``direct'' distinguishes spatial access to a field's fluctuation statistics from inference based only on a secondary macroscopic effect. It does not imply the absence of a preparation model, a measurement basis, optical calibration, or statistical processing. The experiment measures spin-selective atomic densities after controlled operations and infers the field covariance through an explicit relation between those densities and the spin variables.\cite{zhang2026,hung}

Electro-optic studies provide a useful comparison: measurements of electromagnetic vacuum correlations and theoretical descriptions of their sampling likewise specify a coupling between field observables and a detector. The common lesson is not that all vacuum measurements are identical, but that a vacuum observable is always assigned to a specified field, state, mode definition, and readout. Nothing in that usage identifies the vacuum with an absence of physical structure.\cite{benea,moskalenko}

\subsection{Evidence available for the present assessment}
The present work uses the experimental article and its supplementary methods, the related simulator study, and the necessary confinement, Gaussian-state, imaging, and dimensionality literature. It contains no new atomic images and performs no independent regression on the original experimental image stacks. Derived numerical values use published rounded inputs. Prospective tests and illustrative parameter variations are explicitly separated from reported observations. This separation is especially important for the temperature and detection-efficiency estimates.\cite{zhang2026}

\section{Governing TCGS--SEQUENTION architecture}
\label{sec:framework}
\subsection{Source, physical shadow, and effective representation}
The governing architecture distinguishes a complete non-temporal source $(C^4,G,\Psi)$ from the physical shadow carrier $\Sigma^3$. An admissible immersion
\begin{equation}
 X_\sigma:\Sigma^3_\sigma\hookrightarrow C^4,
 \qquad (g_\sigma,\psi_\sigma)=X_\sigma^*(G,\Psi)
 \label{eq:pullback}
\end{equation}
specifies the geometric inheritance of shadow quantities. The direction of $X_\sigma$ must be preserved: it is an immersion from the carrier into the source representation. Its pullback acts in the opposite direction on fields and tensors. Neither operation is, by notation alone, a quotient of source configurations onto measurement outcomes.\cite{foundation}

The framework's four axioms supply the governing roles: complete source content; a conserved singular set $S=\operatorname{Orb}(p_0)$; three-dimensional shadow realization with non-ontic ordering labels; and a single source-grounded constitutive architecture rather than independently added source sectors. These roles are not re-estimated from the BEC measurements. The laboratory Hamiltonian is a domain-level representation within the physical shadow.\cite{foundation,logical}

The source-dimensional argument also has a precise premise structure. A rank-three carrier alone does not yield a four-dimensional source. An independent normal relation is needed to exceed three dimensions, minimal smooth completion selects four rather than a larger dimension, and the exclusion of ontic time classifies that extra source direction as non-temporal. The atomic reduction studied below supplies none of these premises by itself. It instead illustrates a distinct downstream issue: suppression of an observable coordinate need not eliminate its physical carrier dependence.\cite{logical,planck}

\subsection{The typed constitutive relation}
For an admissible complete source configuration $\Xi$, the modern formulation distinguishes
\begin{align}
 M_\Xi&=\mathrm{Info}(\Xi),\nonumber\\
 E_S&=\mathrm{Read}_{S,G}\circ\mathrm{Info},\nonumber\\
 E_{S,\Xi}&=E_S(\Xi),\qquad
 E^\sigma_{S,\Xi}=\mathrm{ev}_\sigma E_{S,\Xi}.
 \label{eq:ecl}
\end{align}
The informational organization, the constitutive map, its complete profile, and one selected readout have different types. With a physical-registration map $\pi^{\mathrm{phys}}_\sigma$, 4D-Counterduction is the composition
\begin{equation}
 \mathrm{Ctd}^{\sigma}_{4,S,G}
 =\pi^{\mathrm{phys}}_\sigma\circ\mathrm{ev}_\sigma
   \circ\mathrm{Read}_{S,G}\circ\mathrm{Info}.
 \label{eq:counterduction}
\end{equation}
This composition is an order of constitutive dependence, not a propagation process through source time.\cite{foundation}

The distinction corrects two possible excesses in a vacuum interpretation. The Extrinsic Constitutive Law (ECL) is not measured as the fluctuation spectrum, and it is not an additional random force acting on the condensate. Conversely, keeping the ECL in its governing role does not mean excluding the experiment from the framework. It means that the spin Hamiltonian, covariance, and imaging operation are domain-level maps, not identities with the complete source-grounded law. No biological or consciousness-specific connection is required for this atomic analysis.\cite{foundation}

\subsection{Three distinct meanings of time}
The framework requires
\begin{equation}
 \mathcal F\not\equiv\lambda_g\not\equiv\topc,
 \qquad \lambda_g\sim h(\lambda_g),\quad h'>0.
 \label{eq:time_types}
\end{equation}
Here $\mathcal F$ is an admissible geometric family of readouts, $\lambda_g$ is a relabelable index, and $\topc$ is an operational clock functional constructed from physical records. The field equations below use laboratory time $t$ in the last sense. Retaining that variable is indispensable for specifying RF pulses, ramps, frequencies, and delays; it does not make the variable a dimension of Counterspace.\cite{foundation,minkowski}

\subsection{What counts as an empirical anchor}
\begin{definition}[Domain-specific empirical anchor]
An empirical anchor is a documented observation together with an explicit, independently assessable bridge to a restricted framework distinction. Its scope is the distinction preserved by that bridge, not every foundational premise of the framework.
\end{definition}
This usage follows the framework's separation between governing necessities, particular mathematical maps, and individual readouts. The anchor here is the coexistence of reduced planar field behaviour with finite transverse carrier structure. A claim of comparative empirical preference over another ontology would additionally require different predicted probabilities for an observation.\cite{foundation,planck}

This last requirement does not weaken the local result. A carefully delimited anchor can be both explanatory and technically useful. It directs attention to the retained dependence on a suppressed coordinate and to the experiment needed to recover that information. It also avoids borrowing the well-tested quantum covariance as evidential authority for a source construction that has not been evaluated by this particular measurement.

\section{Dimensionality must be assigned to a specified object}
\label{sec:types}
The same experiment can legitimately contain objects described as zero-, one-, two-, three-, or infinite-dimensional without asserting that all are autonomous physical substrates. A trap centre is represented by a point; a propagation path by a curve; a planar field by functions of two coordinates; the confined gas by a spatial wavefunction; and the many-body state by a Hilbert-space vector. The meaning of dimension changes with the object and the equivalence relation. The framework's carrier statement concerns physical shadow realization, not the prohibition of lower-dimensional mathematical objects.\cite{foundation,planck}

\begin{longtable}{@{}L{0.19\textwidth}L{0.31\textwidth}L{0.43\textwidth}@{}}
\caption{Dimensional types relevant to vacuum-fluctuation imaging. The rows are not alternative estimates of one common quantity.\label{tab:dimensions}}\\
\toprule Object & Relevant dimension & Permitted interpretation\\\midrule
\endfirsthead
\toprule Object & Relevant dimension & Permitted interpretation\\\midrule
\endhead
\bottomrule\endfoot
Physical atomic carrier & Three spatial directions with anisotropic confinement & A nonzero transverse wavefunction remains even when excited transverse motion is suppressed.\\
Planar spin field & Two in-plane coordinates & An effective field sector obtained after a specified transverse reduction.\\
Density image & Two detector coordinates & A spatial record with a point-spread function and, in the atomic model, a column-density interpretation.\\
Single spin mode & One canonical pair & Two-dimensional phase space; this is not a count of spatial dimensions.\\
Many-body state & Hilbert-space dimension & The number of state-space degrees of freedom, not the dimension of physical space.\\
Spectral estimator & Operator- and scale-dependent $d_s$ & A scaling property of a chosen propagation or diffusion problem.\\
Boundary or defect & Codimension relative to a carrier & A geometric subset does not determine the dimension of the surrounding physical system.\\
Counterspace & Four non-temporal source dimensions within \TCGS & Governing source architecture, distinct from the atomic transverse coordinate and from $2+1$ field notation.\\
\end{longtable}

Three dimensions of the material carrier and two effective coordinates of a field are therefore compatible statements. The experimental literature's use of ``two-dimensional'' is not, in itself, a category error. The error would occur only if the success of that effective description were used to infer that the atoms or the whole apparatus had become an exactly zero-thickness physical surface. The confinement literature explicitly retains the frozen-direction oscillator scale.\cite{hadzibabic,petrov}

Likewise, the \TCGS\ statement that points, lines, and surfaces are mathematical abstractions within the shadow must be read at the level of exact physical realization. It does not deny the empirical utility of a point-defect coordinate, a domain-wall centre, or an ideal interface. A surface can be the exact mathematical level set of an effective field without being an independently existing material sheet of vanishing width. The Planck analysis makes the analogous distinction between a two-dimensional boundary and its three-dimensional admissible domain.\cite{planck}

This distinction also prevents an inversion of the map--territory argument. Calling a coordinate system a map does not render its predictions fictitious. The reduced spin field captures real correlations and experimentally controllable dynamics. The issue is whether the field coordinates exhaust the physical realization, not whether the field theory is useful or empirically meaningful.\cite{zhang2026,foundation}

\section{From a three-dimensional atomic field to a planar effective sector}
\label{sec:reduction}
\subsection{The transverse expansion}
Let $\boldsymbol r=(x,y)$ and let $z$ denote the confined spatial direction. For each spin component $\varsigma$, expand the three-dimensional atomic field in a complete orthonormal transverse basis:
\begin{equation}
 \widehat\Psi_{\varsigma}(\boldsymbol r,z)
 =\sum_{\nu=0}^{\infty}f_\nu(z)
   \widehat\psi_{\varsigma\nu}(\boldsymbol r),
 \qquad \int f_\nu^*(z)f_\mu(z)\,\dd z=\delta_{\nu\mu}.
 \label{eq:mode_expansion}
\end{equation}
For harmonic confinement the lowest transverse eigenfunction is
\begin{equation}
 f_0(z)=\frac{1}{(\pi\lop^2)^{1/4}}
 \exp\!\left(-\frac{z^2}{2\lop^2}\right),
 \qquad \lop=\sqrt{\frac{\hbar}{m\omega_z}}.
 \label{eq:gaussian}
\end{equation}
Retaining the $\nu=0$ sector gives the familiar quasi-two-dimensional reduction. Its validity is a dynamical and energetic approximation, whereas the existence of the nonzero transverse profile is explicit in the construction.\cite{hadzibabic,petrov}

The approximation
\begin{equation}
 \widehat\Psi_{\varsigma}(\boldsymbol r,z)
 \simeq f_0(z)\widehat\psi_{\varsigma}(\boldsymbol r)
 \label{eq:factorization}
\end{equation}
does not set $z$ equal to zero. It assigns the transverse coordinate a specified normalized dependence. In particular,
\begin{equation}
 \langle z\rangle_0=0,\qquad
 \Var_0(z)=\frac{\lop^2}{2}>0
 \quad\text{for finite }m>0,\;0<\omega_z<\infty.
 \label{eq:width}
\end{equation}
The vanishing mean position specifies the symmetry plane, not the thickness. Equating the two would confuse a first moment with a second moment.

The field reduction is not an assertion that all transverse excitations vanish under every experimental operation. It is a choice of retained sector whose accuracy must be assessed for the preparation, propagation, and readout separately. In particular, a large internal-state RF frequency does not by itself excite transverse motion, because the relevant transition also requires a coupling matrix element to a transverse mode. Frequency ratios alone do not establish or refute the reduction.\cite{zhang2026,precursor}

\subsection{An isometric embedding is not a spatial identification}
For a normalized $f\in L^2(\R)$ define
\begin{equation}
 J_f:L^2(\R^2)\longrightarrow L^2(\R^3),
 \qquad (J_f u)(\boldsymbol r,z)=u(\boldsymbol r)f(z).
 \label{eq:isometry}
\end{equation}
Its adjoint is the transverse-mode extraction
\begin{equation}
 (J_f^*v)(\boldsymbol r)=\int f^*(z)v(\boldsymbol r,z)\,\dd z.
 \label{eq:adjoint}
\end{equation}
Direct integration gives $J_f^*J_f=I$ and $P_f:=J_fJ_f^*$ an orthogonal projector onto the selected product sector. These are Hilbert-space statements. They do not assert an isomorphism of the underlying two- and three-dimensional spatial manifolds.

\begin{proposition}[Exact reduced observables on a prescribed sector]
For any bounded operator $A$ on $L^2(\R^2)$ and any normalized $u$, the lifted observable $\widetilde A=J_fAJ_f^*$ satisfies
\begin{equation}
 \langle J_fu,\widetilde A J_fu\rangle
 =\langle u,Au\rangle.
 \label{eq:expectation_lift}
\end{equation}
Thus exact equality of a family of reduced expectation values can coexist with a nonzero transverse profile.
\end{proposition}
\begin{proof}
Insert $J_f^*J_f=I$ on both sides of $A$. The conclusion holds for the specified subspace and lifted observables; it does not state that arbitrary three-dimensional observables or dynamics are exhausted by this subspace.
\end{proof}

The proposition is a simple representation result, not a new physical law. Its relevance is inferential: even exact success of a planar observable algebra would not imply that the corresponding three-dimensional realization had lost its transverse direction. The missing information resides in the choice of $f$ and in observables outside the retained algebra.

\subsection{The effective interaction retains transverse information}
For weak contact interactions, the three-dimensional coupling is $g_{3D}=4\pi\hbar^2a/m$, with the appropriate scattering length $a$. Substitution of Equation~\eqref{eq:factorization} into the contact-interaction energy yields
\begin{align}
 g_{2D}&=g_{3D}\int_{\R}|f_0(z)|^4\,\dd z
       =\frac{g_{3D}}{\sqrt{2\pi}\lop}\nonumber\\
       &=\sqrt{8\pi}\,\frac{\hbar^2a}{m\lop}.
 \label{eq:g2d}
\end{align}
This is the coupling form used in the experimental supplementary material. It is the weak quasi-two-dimensional reduction, not a universal energy-independent scattering formula in every strongly confined regime.\cite{zhang2026,hadzibabic,petrov}

The dimensional check is decisive. The three-dimensional coupling has units of energy times volume, the overlap integral has units of inverse length, and the reduced coupling has units of energy times area. The eliminated coordinate remains in a parameter with measurable consequences. It has been integrated out of the independent coordinate list, not erased from the physical realization.

For the two-component gas, the same reduction applies componentwise to $g_{\uparrow\uparrow}$, $g_{\uparrow\downarrow}$, and $g_{\downarrow\downarrow}$. The interaction energy can then be written
\begin{equation}
 \mathcal E_{\mathrm{int}}
 =\frac{g_dn^2}{2}+\frac{\kappa n^2}{2}(Z-Z_0)^2,
 \label{eq:interaction}
\end{equation}
where
\begin{align}
 \kappa&=\frac{g_{\uparrow\uparrow}+g_{\downarrow\downarrow}
                   -2g_{\uparrow\downarrow}}{4},\nonumber\\
 g_d&=\frac{g_{\uparrow\uparrow}g_{\downarrow\downarrow}
                   -g_{\uparrow\downarrow}^2}
                  {g_{\uparrow\uparrow}+g_{\downarrow\downarrow}
                   -2g_{\uparrow\downarrow}}.
 \label{eq:kappa}
\end{align}
At the operating imbalance, perturbative spin and density dynamics decouple to the order considered by the authors. Both interaction scales nevertheless inherit the transverse overlap.\cite{zhang2026}

\subsection{Two reductions that must not be confused}
The transverse-mode extraction $J_f^*$ acts on a field amplitude. Absorption imaging instead registers a density-related signal, with a column-density interpretation in the atomic model:
\begin{equation}
 n_{2D}(\boldsymbol r)
 =\int n_{3D}(\boldsymbol r,z)\,\dd z.
 \label{eq:column}
\end{equation}
For a normalized product profile, $n_{3D}=|f_0|^2n_{2D}$. The equality is not obtained by evaluating $n_{3D}$ at the plane $z=0$; that evaluation would have different units and normalization. The microscope adds its own optical response and calibration to this density relation.\cite{zhang2026,hung}

These two operations occur inside the physical shadow. They are neither the immersion $X_\sigma$ in Equation~\eqref{eq:pullback} nor the full constitutive relation in Equation~\eqref{eq:counterduction}. Similar arrows in a diagram cannot substitute for matching the domains and codomains of the actual maps.\cite{foundation}

\section{Finite-resource carrier irreducibility}
\label{sec:irreducibility}
\subsection{A precise local meaning of irreducibility}
\begin{definition}[Finite-resource transverse irreducibility]
A family of massive-atom states is transversely irreducible at finite kinetic resource if no sequence in that family can approach zero transverse variance while maintaining a uniform finite bound on its transverse kinetic energy.
\end{definition}
The definition is deliberately local to a specified physical model. It does not identify irreducibility with the impossibility of writing a lower-dimensional theory, and it does not assert a universal minimum length. It identifies what a zero-width limit would cost in the quantum-mechanical carrier used by the experiment.

\begin{proposition}[Finite-energy width bound]
Let a normalized atomic state have a well-defined transverse position variance $\sigma_z^2>0$ and finite transverse kinetic energy
\begin{equation}
 K_z=\frac{\langle p_z^2\rangle}{2m},\qquad m>0.
 \label{eq:kinetic}
\end{equation}
Under the canonical position--momentum uncertainty relation,
\begin{equation}
 K_z\geq\frac{\hbar^2}{8m\sigma_z^2}.
 \label{eq:energy_bound}
\end{equation}
Consequently, a uniform bound $K_z\leq K_{\max}<\infty$ implies
\begin{equation}
 \sigma_z\geq\frac{\hbar}{\sqrt{8mK_{\max}}}>0.
 \label{eq:minimum_width}
\end{equation}
\end{proposition}
\begin{proof}
The uncertainty relation gives $\Var(z)\Var(p_z)\geq\hbar^2/4$. Since $\langle p_z^2\rangle\geq\Var(p_z)$, division by $2m$ yields Equation~\eqref{eq:energy_bound}. Rearranging under the assumed upper bound gives Equation~\eqref{eq:minimum_width}. A nonzero mean momentum strengthens the energy requirement rather than evading it.
\end{proof}

The Gaussian harmonic ground state saturates the bound when its mean momentum vanishes. Its kinetic and potential energies are each $\hbar\omega_z/4$, and its variance is Equation~\eqref{eq:width}. The connection between confinement, width, and momentum uncertainty is therefore exact for the nominal harmonic state, although the experimental realization may depart from that ideal profile. The uncertainty structure is the same canonical structure retained in the covariance analysis below.\cite{hadzibabic,weedbrook}

\subsection{A zero-thickness surface is not an ordinary volume-normalized state}
\begin{proposition}[Measure-zero support obstruction]
Let $S_2\subset\R^3$ be a regular two-dimensional surface, hence a set of three-dimensional Lebesgue measure zero. A function $u\in L^2(\R^3,\dd^3x)$ supported on $S_2$ is the zero vector of that Hilbert space and cannot be a normalized atomic wavefunction.
\end{proposition}
\begin{proof}
Because $u=0$ almost everywhere outside $S_2$ and $S_2$ has zero volume measure,
$\|u\|_2^2=\int_{S_2}|u|^2\dd^3x=0$.
\end{proof}

An intrinsic surface theory on $L^2(S_2,\dd A)$ is mathematically legitimate. It uses a different measure and a separately defined operator domain. The proposition does not prohibit such theories. It shows why an intrinsic surface state must not be identified with an unchanged volume-normalized atomic state whose width has somehow become exactly zero.

Similarly, a distribution such as $\delta(z)$ is not a normalized element of $L^2(\R,\dd z)$. A sequence of transverse probability densities can converge weakly to a delta measure, but the corresponding amplitudes need not converge to a normalizable zero-width vector in the original Hilbert space. Weak convergence of probability measures is not preservation of the original kinetic-energy domain.

\subsection{The singular limit and its cost}
For a normalized profile $f$, define $f_\epsilon(z)=\epsilon^{-1/2}f(z/\epsilon)$. When the moments and kinetic form exist,
\begin{equation}
 \sigma_z^2[f_\epsilon]=\epsilon^2\sigma_z^2[f],
 \qquad K_z[f_\epsilon]=\epsilon^{-2}K_z[f].
 \label{eq:scaling_width}
\end{equation}
The zero-width limit is therefore an unbounded-energy limit in this model. For harmonic confinement, $\lop\to0$ at fixed mass corresponds to $\omega_z\to\infty$, which likewise makes the transverse zero-point energy diverge.

Subtracting the transverse ground-state energy when defining an effective planar Hamiltonian is a useful energy-reference convention. It does not remove the physical confining field, change the wavefunction width, or undo the momentum spread. A subtraction in the reduced energy origin and a removal of physical localization cost are different operations.

The scope is also important at very high energy. The nonrelativistic atomic description ceases to be adequate before an actual laboratory procedure reaches arbitrarily small lengths. The present result is a bound within the stated model, not an extrapolation through atomic breakup, relativistic particle production, or quantum gravity. No Planck-scale lower bound or theorem about all conceivable lower-dimensional fundamental theories is inferred.\cite{planck}

\subsection{The anchor supported by these results}
The experimental carrier has finite confinement, and its effective coupling explicitly contains the confinement length. Equations~\eqref{eq:width}, \eqref{eq:g2d}, and \eqref{eq:energy_bound} therefore supply a precise bridge from the reported apparatus to the carrier--representation distinction. Under \TCGS, this is the relevant finite-resource realization of three-dimensional carrier irreducibility: the planar description does not convert its material support into an autonomous two-dimensional substrate.\cite{zhang2026,foundation}

This conclusion is narrower than a universal ontological prohibition, but stronger than a purely verbal observation about the word ``planar.'' It yields a nonzero nominal width, an energy bound, a confinement-dependent interaction, and an experimentally accessible transverse response. Each quantity belongs to a defined model and can be checked independently of the source interpretation.

\section{Quantitative audit of the transverse carrier}
\label{sec:numerics}
\subsection{Nominal scales and precision}
Using $m\simeq39u$, with $u\simeq1.66054\times10^{-27}\,\mathrm{kg}$, and the reported $\omega_z=2\pi\times10^3\,\mathrm{s}^{-1}$ gives
\begin{align}
 \lop&\simeq5.09\times10^{-7}\,\mathrm m
          =0.509\,\mu\mathrm m,\nonumber\\
 \sigma_z&=\lop/\sqrt2\simeq0.360\,\mu\mathrm m,\nonumber\\
 \frac{\hbar\omega_z}{k_B}&\simeq48.0\,\mathrm{nK}.
 \label{eq:numeric_width}
\end{align}
These are calculated consequences of a nominal harmonic model, not a direct experimental measurement of thickness. The isotope mass approximation is far smaller than the uncertainty implied by a trap frequency reported only as approximately $1\,\mathrm{kHz}$.\cite{zhang2026}

With $L\simeq33\,\mu\mathrm m$, the aspect ratios are $\lop/L\simeq0.0154$ and $\sigma_z/L\simeq0.0109$. Strong anisotropy is therefore entirely compatible with a nonzero third-direction profile. The harmonic ground-state energy corresponds to $24.0\,\mathrm{nK}$ and its kinetic contribution to $12.0\,\mathrm{nK}$ when each is divided by $k_B$. The latter quantity is not the gas temperature: a zero-temperature quantum state can have nonzero kinetic energy.\cite{zhang2026}

\begin{table}[htbp]
\centering\small
\caption{Nominal carrier audit. Reported values are distinguished from calculations under the harmonic, lowest-transverse-mode approximation.}
\label{tab:scales}
\begin{tabular}{@{}L{0.40\textwidth}L{0.22\textwidth}L{0.29\textwidth}@{}}
\toprule Quantity & Value & Status\\\midrule
Transverse frequency $\omega_z/(2\pi)$ & $\simeq1\,\mathrm{kHz}$ & Reported\cite{zhang2026}\\
In-plane side $L$ & $\simeq33\,\mu\mathrm m$ & Reported\cite{zhang2026}\\
Oscillator length $\lop$ & $0.509\,\mu\mathrm m$ & Derived nominal value\\
RMS transverse width $\sigma_z$ & $0.360\,\mu\mathrm m$ & Derived nominal value\\
Axial excitation gap divided by $k_B$ & $48.0\,\mathrm{nK}$ & Derived nominal value\\
Axial zero-point energy divided by $k_B$ & $24.0\,\mathrm{nK}$ & Derived nominal value\\
Axial kinetic energy divided by $k_B$ & $12.0\,\mathrm{nK}$ & Derived nominal value\\
$\mu_d/(\hbar\omega_z)$ & $\simeq0.120$ & Ratio of reported scales\\
$\mu_s/(\hbar\omega_z)$ & $\simeq0.900$ & Ratio of reported scales\\
$\musp/(\hbar\omega_z)$, $\alpha=0.72$ & $\simeq0.648$ & Scaled spin ratio\\
Spin-sector temperature fit & $(6\pm6)\,\mathrm{nK}$ & Source fit, not refitted\\
\bottomrule
\end{tabular}
\end{table}

\subsection{A dimensional hierarchy must not be overstated}
The density scale is appreciably below the transverse gap, but the unscaled spin-interaction scale is approximately $0.9$ times that gap. The scaled spin quantity is approximately $0.648$ times the gap. Consequently, the source does not justify a blanket assertion that \emph{every} spin, interaction, and readout scale is asymptotically smaller than $\hbar\omega_z$. The common profile and the decoupling of perturbative density and spin fluctuations must be evaluated with the actual coupling structure.\cite{zhang2026,precursor}

This caveat changes neither the nonzero width nor the distinction between the effective field and its carrier. It changes the strength of the quantitative approximation. A controlled reduction needs to estimate population and virtual admixture in higher transverse sectors, not merely attach a ``2D'' label to the apparatus. The projected Hamiltonian construction in Appendix~\ref{app:projection} identifies the additional hypotheses required for an operator-level error estimate.

The temperature inferred from spin-mode amplitudes also need not certify a common equilibrium temperature for every external and internal degree of freedom. The experimental fit constrains the effective spin occupations under the model used in that fit. It cannot, without a thermalization assumption, be converted into a measured upper bound on all transverse occupations.\cite{zhang2026}

\subsection{Sensitivity to confinement uncertainty}
At fixed nominal mass,
\begin{equation}
 \frac{\dd\lop}{\lop}=-\frac12\frac{\dd\omega_z}{\omega_z},
 \qquad
 \frac{\dd g_{2D}}{g_{2D}}
 =\frac{\dd a}{a}+\frac12\frac{\dd\omega_z}{\omega_z}.
 \label{eq:sensitivity}
\end{equation}
Thus a hypothetical 10\% frequency uncertainty would produce approximately 5\% width uncertainty to first order. This is a sensitivity illustration, not an uncertainty assigned to the authors' apparatus. The paper does not supply the trap-frequency covariance needed to attach a defensible metrological error bar to Equation~\eqref{eq:numeric_width}.

For an uncertainty model including mass, scattering length, and frequency, the covariance terms must be propagated rather than assuming independent calibration. The nominal width is useful because it demonstrates finite support scale by a large margin; it should not be presented with excessive significant figures or as a newly measured condensate thickness.

\section{The spin field as a family of canonical oscillators}
\label{sec:spin}
\subsection{Scaled variables and equations of motion}
The interaction-minimizing imbalance in the experiment is $Z_0\simeq0.53$. Following the source, define
\begin{equation}
 \alpha=1-Z_0^2\simeq0.72,\qquad
 \delta Z'=\frac{\delta Z}{\alpha},\qquad
 \Omega'=\frac{\Omega}{\sqrt\alpha},\qquad
 \musp=\alpha\mu_s.
 \label{eq:scaled}
\end{equation}
For a planar wavevector $\boldsymbol k$, let
\begin{equation}
 \epsilon_k=\frac{\hbar^2k^2}{2m},\qquad
 a_k=\epsilon_k+\hbar\Omega',\qquad
 b_k=a_k+2\musp.
 \label{eq:ab}
\end{equation}
The linear spin equations are
\begin{equation}
 \hbar\dot{\delta\phi}_k=b_k\delta Z'_k,
 \qquad
 \hbar\dot{\delta Z'}_k=-a_k\delta\phi_k.
 \label{eq:spin_motion}
\end{equation}
These are the authors' perturbative equations, with $a_k$ and $b_k$ introduced only to simplify the derivation.\cite{zhang2026}

For positive $a_k$ and $b_k$, differentiation gives
\begin{equation}
 \ddot{\delta\phi}_k+\omega_k^2\delta\phi_k=0,
 \qquad \hbar\omega_k=\sqrt{a_kb_k},\qquad
 r_k=\sqrt{b_k/a_k}.
 \label{eq:dispersion}
\end{equation}
The exactly gapless uniform mode at $k=0$, $\Omega'=0$ is excluded from the normalizable harmonic-ground-state calculation. Its compact phase and finite-size treatment require separate care. A divergence in the corresponding harmonic formula is a limit of that chart, not evidence for a zero-dimensional or singular physical particle.

\subsection{Real modes and canonical normalization}
The experimental Fourier convention assigns the real-space field a complex coefficient for each $\boldsymbol k$, with the usual relation between $\boldsymbol k$ and $-\boldsymbol k$. For a canonical covariance calculation, choose one orthonormal real standing-wave mode at a time. If $\phi_j$ and $z_j$ denote its relative-phase and scaled-imbalance amplitudes, the corresponding normalization is
\begin{equation}
 [\phi_j,z_{j'}]=\frac{2\ii}{\alpha\nb}\delta_{jj'}.
 \label{eq:commutator}
\end{equation}
This avoids treating the complex Fourier coefficient and its conjugate as two independent Hermitian observables. It is a change of mode basis, not a change of experimental content.\cite{zhang2026}

Define the dimensionless canonical pair
\begin{equation}
 Q_j=\sqrt{\frac{\alpha\nb}{2}}\phi_j,
 \qquad P_j=\sqrt{\frac{\alpha\nb}{2}}z_j,
 \qquad [Q_j,P_{j'}]=\ii\delta_{jj'}.
 \label{eq:QP}
\end{equation}
Then the quadratic mode Hamiltonian is
\begin{equation}
 H_j=\frac12\left(a_jQ_j^2+b_jP_j^2\right),
 \label{eq:Hmode}
\end{equation}
which reproduces Equation~\eqref{eq:spin_motion}. Its use of dimensionless quadratures explains why $a_j$ and $b_j$ have energy units and why $\hbar\omega_j=\sqrt{a_jb_j}$.

\subsection{Ground-state and thermal covariance}
With $r=\sqrt{b/a}$, the canonical rescaling $q=Q/\sqrt r$, $p=\sqrt r P$ gives
\begin{equation}
 H=\frac{\sqrt{ab}}{2}(q^2+p^2).
 \label{eq:isotropicH}
\end{equation}
For a centred thermal oscillator state with occupation $N$, its symmetrized covariance in the original quadratures is therefore
\begin{equation}
 C_k=\left(N_k+\frac12\right)
 \begin{pmatrix}r_k&0\\0&r_k^{-1}\end{pmatrix}.
 \label{eq:covariance}
\end{equation}
The vacuum is $N_k=0$. This is a canonical Gaussian-state calculation; its physical interpretation is fixed by the mode Hamiltonian and by the chosen quadrature normalization.\cite{weedbrook}

Transforming back to the experimental field amplitudes gives
\begin{align}
 V_{\phi,k}&:=\langle|\delta\phi_k|^2\rangle
       =\frac{2N_k+1}{\alpha\nb}\,r_k,\nonumber\\
 V_{Z',k}&:=\langle|\delta Z'_k|^2\rangle
       =\frac{2N_k+1}{\alpha\nb}\,r_k^{-1}.
 \label{eq:field_variances}
\end{align}
At $N_k=0$ these reproduce Equation~(4) of the experiment. The inverse density factor follows the source's mode normalization. It should not be removed when comparing an areal field amplitude with a canonically normalized oscillator.\cite{zhang2026}

\subsection{Stationarity is compatible with nonzero variance}
For a ground-state density operator $\rho_0$ of a fixed Hamiltonian, $[H,\rho_0]=0$, so the state is stationary under its own evolution. Nevertheless, $\Tr(\rho_0Q^2)$ and $\Tr(\rho_0P^2)$ are nonzero. A stationary distribution of outcomes is not the same as a classical trajectory of a field value that remains identically zero. The atomic experiment accesses those distributions through repeated measurements and controlled changes of Hamiltonian or quadrature.\cite{zhang2026,weedbrook}

This is the precise level at which \TCGS\ interprets quantum uncertainty as a constraint on accessible shadow statistics. The interpretation does not delete the commutator, the vacuum variance, or the experimentally verified response. It separates those operational statements from a claim about the source itself undergoing random temporal motion.\cite{foundation,born}

\section{What the covariance measurement preserves}
\label{sec:covariance}
\subsection{Canonical invariants and quadrature-dependent quantities}
For a real canonical pair $R=(Q,P)^T$, define the centred symmetrized covariance by
\begin{equation}
 C_{ab}=\frac12\langle\Delta R_a\Delta R_b+\Delta R_b\Delta R_a\rangle.
 \label{eq:Cdefinition}
\end{equation}
With $[Q,P]=\ii$, the uncertainty condition gives $\det C\geq1/4$. Equation~\eqref{eq:covariance} gives
\begin{equation}
 \det C_k=\left(N_k+\frac12\right)^2.
 \label{eq:detC}
\end{equation}
The ground-state determinant is $1/4$, although the two diagonal entries need not equal one another. This convention differs by a factor of two in quadrature normalization from conventions in which the vacuum covariance is the identity.\cite{weedbrook}

A real symplectic transformation $S$ satisfies $SJS^T=J$, where
\begin{equation}
 J=\begin{pmatrix}0&1\\-1&0\end{pmatrix},\qquad C\longmapsto SCS^T.
 \label{eq:symplectic}
\end{equation}
For one mode, $\det S=1$, so $\det C$ is unchanged. By contrast, $C_{QQ}$, $C_{PP}$, and their difference generally change. Under an orthogonal phase-space rotation, the trace also stays fixed; under a general squeeze, it need not. The admissible transformation class must therefore be stated before a quantity is called invariant.\cite{weedbrook}

This elementary distinction has a direct \TCGS\ implication. A repeatable spectrum within one protocol is not automatically a source-level invariant under arbitrary selectors. Some stability follows from canonical mathematics, some from fixing the Hamiltonian and state, and some must be demonstrated across physically equivalent readouts. The framework requires those classes to remain distinct.\cite{foundation,planck}

\subsection{Quadrature rotation and the observed amplitude}
For an ideal readout rotation, let $e_\vartheta=(\cos\vartheta,\sin\vartheta)^T$. The accessible variance is
\begin{equation}
 V(\vartheta)=e_\vartheta^TCe_\vartheta
 =\frac{\Tr C}{2}
 +\frac{C_{QQ}-C_{PP}}{2}\cos2\vartheta
 +C_{QP}\sin2\vartheta.
 \label{eq:tomography}
\end{equation}
Its peak-to-peak amplitude is
\begin{equation}
 A_{\mathrm{quad}}
 =\sqrt{(C_{QQ}-C_{PP})^2+4C_{QP}^2}.
 \label{eq:general_amplitude}
\end{equation}
For the equilibrium covariance in Equation~\eqref{eq:covariance}, $C_{QP}=0$ and $r_k\geq1$, giving $A_{\mathrm{quad}}=(N_k+1/2)(r_k-r_k^{-1})$.

The experimentally defined structure factor satisfies
\begin{equation}
 S_Y(k,t)=\frac{S_d(k)+\alpha\nb V_{\phi,k}(t)}{2},
 \label{eq:SY}
\end{equation}
where the decoupled total-density contribution $S_d$ is independent of the readout delay in the model. The spin contribution equals the normalized $Q$ variance. Since
\begin{equation}
 r_k-r_k^{-1}=\frac{b_k-a_k}{\sqrt{a_kb_k}}
 =\frac{2\musp}{\hbar\omega_k},
 \label{eq:anisotropy_identity}
\end{equation}
the ideal equilibrium readout amplitude is
\begin{equation}
 A_Y(k)=\frac{2\musp}{\hbar\omega_k}\left(N_k+\frac12\right).
 \label{eq:AY}
\end{equation}
This reproduces the source's direct-readout formula without identifying the amplitude with a full covariance determinant.\cite{zhang2026}

\subsection{Anisotropy is not a complete state diagnostic}
Measuring $A_Y$ robustly avoids a time-independent density offset, which is a major experimental advantage. However, a difference of variances does not determine both variances separately. For example, covariances
\begin{equation}
 C(c)=\begin{pmatrix}c+A&0\\0&c\end{pmatrix},\qquad c>0,
 \label{eq:ambiguity_cov}
\end{equation}
all have the same positive quadrature amplitude $A$, while their determinants $c(c+A)$ differ. Restricting to values satisfying the uncertainty bound still leaves a family of distinct states. Thus the amplitude identifies a thermal occupation only after the Hamiltonian-dependent equilibrium covariance form has been assumed.

The reported agreement with a vacuum-dominated thermal model is accordingly substantial but model-specific. It is not a model-independent reconstruction of every covariance entry or an independent purity measurement. Full quadrature tomography would require an absolute variance calibration and enough independent quadratures to determine $C_{QQ}$, $C_{PP}$, and $C_{QP}$. For a Gaussian single-mode state in the present normalization, purity would then be $(2\sqrt{\det C})^{-1}$.\cite{zhang2026,weedbrook}

\subsection{Finite readout coupling}
A strong readout coupling makes the phase-space evolution nearly circular, but ``nearly'' is not the same as exactly. For a finite readout Hamiltonian with ratio $r_m=\sqrt{b_m/a_m}$, the propagator is
\begin{equation}
 S_m(t)=\begin{pmatrix}
 \cos\omega_mt&r_m\sin\omega_mt\\
 -r_m^{-1}\sin\omega_mt&\cos\omega_mt
 \end{pmatrix}.
 \label{eq:readoutS}
\end{equation}
The peak-to-peak $Q$-variance signal is therefore
\begin{equation}
 A_m=\sqrt{(C_{QQ}-r_m^2C_{PP})^2+4r_m^2C_{QP}^2}.
 \label{eq:finite_rotation}
\end{equation}
Only in the appropriate large-coupling limit does this reduce to Equation~\eqref{eq:general_amplitude}.

For the reported readout scale $\hbar\Omega'_m/\musp\simeq20$, the zero-kinetic-energy estimate gives $r_m\simeq\sqrt{1.1}\simeq1.049$. This is a useful scale for auditing corrections, not a newly measured systematic error. It shows why pulse duration, finite coupling, and actual readout dynamics belong in the forward model. The authors themselves include finite-coupling and rotation corrections in their calibration analysis.\cite{zhang2026}

\section{Vacuum scaling, thermal scaling, and spatial dimension}
\label{sec:spectrum}
\subsection{The exact thermal interpolation}
In thermal equilibrium,
\begin{equation}
 N_k=\frac{1}{e^{x_k}-1},\qquad
 x_k=\frac{\hbar\omega_k}{k_BT}.
 \label{eq:occupation}
\end{equation}
Equation~\eqref{eq:AY} becomes
\begin{equation}
 A_Y(k)=\frac{\musp}{\hbar\omega_k}
 \coth\!\left(\frac{x_k}{2}\right).
 \label{eq:coth}
\end{equation}
For $x_k\gg1$ this gives $A_Y\propto\omega_k^{-1}$, while for $x_k\ll1$ it gives $A_Y\simeq2\musp k_BT/(\hbar^2\omega_k^2)$. These are the two limits distinguished in the experimental paper.\cite{zhang2026}

A useful derived local slope is
\begin{equation}
 \beta_{\mathrm{eff}}
 :=-\frac{\dd\log A_Y}{\dd\log\omega_k}
 =1+\frac{x_k}{\sinh x_k},
 \label{eq:slope}
\end{equation}
when $\musp$ and $T$ are held fixed. It interpolates continuously from two to one. An experiment probing only a narrow frequency range should therefore compare the thermal interpolation with the data rather than treating every departure from an exact integer slope as a different vacuum ontology.

The source reports that the central fitted temperature corresponds to $N_k\lesssim0.1$ outside the nonadiabatic region. At that occupation, the thermal multiplier $2N_k+1$ is at most approximately $1.2$. This is compatible with vacuum dominance, but it is not the statement that every retained mode has exactly zero occupation. The finite uncertainty on $T$ and its physical lower bound at zero should also be retained when interpreting the fit.\cite{zhang2026}

\subsection{Why the exponent is not a spatial-dimension measurement}
Equation~\eqref{eq:coth} is a per-mode covariance result. Its inverse-frequency vacuum limit follows from the oscillator stiffnesses and their difference. It does not count the number of spatial coordinates. An analogous canonical mode in another spatial dimension can have the same covariance-frequency dependence while the number of modes in a shell is different.

For an isotropic scalar-mode problem in $d$ effective spatial dimensions, a local variance involves the mode measure:
\begin{equation}
 \langle\phi^2\rangle
 =\frac{S_{d-1}}{(2\pi)^d}
   \int_0^\infty k^{d-1}V_\phi(k)\,\dd k,
 \label{eq:mode_measure}
\end{equation}
subject to the physical cutoffs and normalization. The dimensional information enters through the mode density, dispersion, boundary conditions, and observable definition. Consequently, the isolated law $A_Y\propto\omega^{-1}$ cannot distinguish a two-dimensional carrier from a three-dimensional carrier with a reduced planar sector.

As an example within a relativistic long-wavelength sector, if $V_\phi(k)\propto1/\omega_k$ and $\omega_k\propto k$, the radial integrand scales as $k^{d-2}$. The same per-mode exponent thus produces different integrated infrared behaviour in different dimensions. Conversely, a change in observable normalization can change the per-mode amplitude without changing the underlying carrier dimension. This is the reason for separating covariance scaling from a genuine dimensional estimator.

\subsection{What the sine-Gordon correspondence retains}
Expanding the measured Bogoliubov dispersion gives
\begin{equation}
 \omega_k^2
 =\Omega'\!\left(\Omega'+\frac{2\musp}{\hbar}\right)
 +\frac{\musp+\hbar\Omega'}{m}k^2
 +\frac{\hbar^2}{4m^2}k^4.
 \label{eq:expanded_dispersion}
\end{equation}
In the Josephson, long-wavelength regime, $\hbar\Omega'\ll\musp$ and $\epsilon_k\ll\musp$, the leading expression is
\begin{equation}
 \omega_k^2\simeq\omega_0^2+c_s^2k^2,
 \qquad \omega_0^2\simeq\frac{2\musp\Omega'}{\hbar},
 \qquad c_s^2\simeq\frac{\musp}{m}.
 \label{eq:relativistic}
\end{equation}
The nonlinear relative-phase potential supplies the sine-Gordon extension beyond the quadratic regime. The correction terms in Equation~\eqref{eq:expanded_dispersion} state where the relativistic approximation must be controlled.\cite{zhang2026,precursor,son}

The effective speed $c_s$ is a property of the atomic spin sector. It is not the electromagnetic invariant speed, and the tunable gap is not an independent determination of a fundamental particle mass. The notation $2+1$ records two effective spatial variables plus a time parameter in the simulator's field equation. It does not identify the simulator with a complete relativistic universe. Analogue models are scientifically valuable precisely because a specified subset of equations and observables can transfer between distinct physical realizations.\cite{analogue}

\section{Metrological audit of the images and structure factors}
\label{sec:metrology}
\subsection{A measurement chain rather than a bare photograph}
The experiment extracts fluctuations from independent optical-density images after subtracting the mean profile. Its methods report Gaussian-apodized windowing, principal-component-based fringe suppression, masking of residual contaminated wavevectors, subtraction of photon shot noise, azimuthal averaging, and normalization to an atomic shot-noise reference. Displayed images are additionally low-pass filtered. These steps should be retained in any interpretation of the visible patterns and their Fourier amplitudes.\cite{zhang2026}

A useful schematic forward chain is
\begin{equation}
 \rho_{3D}\xrightarrow{\mathcal R_f}\rho_{\mathrm{spin}}
 \xrightarrow{\mathcal U_\vartheta}\rho_\vartheta
 \xrightarrow{\mathcal M}\mathrm{OD}_j
 \xrightarrow{\mathcal W,\mathcal F,\mathcal N}\widehat S_Y.
 \label{eq:measurement_chain}
\end{equation}
Here $\mathcal R_f$ denotes a declared state or mode reduction, $\mathcal U_\vartheta$ the actual RF operation, $\mathcal M$ the optical measurement, and the last arrow windowing, Fourier analysis, and normalization. The notation is schematic because the operators depend on the apparatus; it is not a newly calibrated imaging model.

The optical transfer function and the analysis window play different roles. For a homogeneous field and a linear optical response, an ideal power spectrum is multiplied by the squared optical transfer function. Multiplication by a finite real-space window instead mixes Fourier modes through convolution. Their effects cannot generally be represented by one constant amplitude factor. Density-correlation imaging theory provides the relevant technical comparison.\cite{hung}

\subsection{Independent detection loss does not cancel in normalization}
The source decomposes a normalized structure factor into a unit shot-noise term and a connected correlation term,
\begin{equation}
 S(k)=1+S_c(k).
 \label{eq:structure_split}
\end{equation}
If each atom is detected independently with probability $\gamma$, the detected mean number is proportional to $\gamma$, whereas a detected pair is proportional to $\gamma^2$. The normalized connected contribution is therefore attenuated once:
\begin{equation}
 S_{\mathrm{obs}}(k)=1+\gamma S_c(k)
 =1+\gamma[S(k)-1].
 \label{eq:thinning}
\end{equation}
The additional binomial noise preserves the unit shot-noise baseline. Normalizing by an uncorrelated reference does not undo the attenuation of correlations.\cite{zhang2026}

The authors find an approximately $\eta=0.7$ residual amplitude factor relative to their corrected theoretical prediction and identify imperfect atom detection as a likely explanation. This factor is then transferred to the direct-readout comparison. It should be described as a reported calibration factor with a proposed physical explanation, not as a uniquely established detection probability or a discrepancy demanding new vacuum physics.\cite{zhang2026}

For the amplified protocol, the source also reports a separate approximately 20\% reduction associated with finite preparation and rotation parameters, calculated using a truncated-Wigner treatment. That correction and the residual $\eta$ factor refer to different parts of the forward calculation. Their provenance must not be merged into a single unexplained normalization, nor should the first be described as a thermal contribution.\cite{zhang2026,blakie}

\subsection{Adiabaticity and selection of modes}
The direct-readout analysis marks a region where the stated criterion $|\dot\omega_k|<\omega_k^2$ fails. The fit is reported for points outside that region. This is a defensible way to expose the intended preparation domain, but a strict adiabatic theorem generally requires more than the inequality being less than unity: the full ramp, relevant matrix elements, endpoints, and integrated excitation error matter. No new bound on the actual experimental excitation probability is calculated here.\cite{zhang2026}

For the reduced oscillator, define a local ramp parameter
\begin{equation}
 \varepsilon_{\mathrm{ad},k}(t)
 =\frac{|\dot\omega_k(t)|}{\omega_k(t)^2}.
 \label{eq:adiabatic}
\end{equation}
Near a closing gap this parameter grows, even for a slow laboratory ramp. An apparent excess fluctuation there can therefore be preparation-dependent rather than a change in vacuum ontology. The appropriate comparison is between a specified nonadiabatic forward model and the measured covariance, not between the excluded region and an unrestricted zero-temperature formula.

\subsection{Finite ensembles and displayed patterns}
For independent centred Gaussian samples of one real mode, an unbiased sample variance has relative standard error $\sqrt{2/(M-1)}$. At $M=70$ and $120$ this is approximately 17\% and 13\%, respectively. These calculations illustrate finite-ensemble scaling only. The published amplitudes are obtained from time-dependent, azimuthally averaged structure factors and fitted oscillations, so these numbers are not substitutes for the paper's error bars.\cite{zhang2026}

Subtracting a mean estimated from the same $M$ samples also changes the divisor required for an unbiased variance. A divisor $M$ would produce the familiar $(M-1)/M$ expectation factor, whereas $M-1$ removes it in the elementary independent-sample case. The present work does not claim that the authors omitted such a correction; it identifies a reproducibility item that should be visible in any reanalysis of the original image stacks.

The map-level artifact category must therefore be used carefully. A filtered image is a real record processed under a specified rule; it is not necessarily a false observation. A fringe artifact, a window-induced mode mixture, and an actual physical anisotropy have different causes. Calling them all ``projection'' would erase the distinction that the framework is intended to preserve.\cite{foundation,zhang2026}

\section{Operational time without dimensional reclassification}
\label{sec:time}
\subsection{A reparameterization must transform the equations}
Consider a mode equation written in laboratory time,
\begin{equation}
 \frac{\dd R}{\dd t}=A(t)R.
 \label{eq:lab_dynamics}
\end{equation}
If a monotone label $s$ is introduced through $t=t(s)$, the same path is represented by
\begin{equation}
 \frac{\dd R}{\dd s}=\frac{\dd t}{\dd s}A(t(s))R.
 \label{eq:reparameterized}
\end{equation}
The conversion factor is essential. Replacing $t$ by another symbol while leaving a numerical frequency unchanged is not a gauge-covariant rewriting of the same operational dynamics.

For an isolated oscillator with slowly varying or prescribed instantaneous frequency, the accumulated phase is
\begin{equation}
 \Theta=\int_{t_1}^{t_2}\omega(t)\,\dd t
 =\int_{s_1}^{s_2}\omega(t(s))\frac{\dd t}{\dd s}\,\dd s.
 \label{eq:phase_invariant}
\end{equation}
It is this dimensionless phase, together with the actual protocol and observables, that survives a consistent reparameterization. Laboratory delays retain their clock-calibrated physical content even though the source ontology does not assign them an independent temporal direction.\cite{foundation,minkowski}

\subsection{A physical quench is not a gauge transformation}
Changing the RF amplitude changes the physical Hamiltonian of the selected spin sector. Changing a readout pulse changes which observable is registered. Neither intervention should be classified as a mere relabeling of the foliation parameter. The same source-grounded architecture may encompass these different experimental situations, but their observable states and covariance matrices can genuinely differ.\cite{zhang2026,foundation}

This distinction is necessary for the invariant analysis. A determinant preserved by an ideal symplectic unitary is preserved because of the mathematical structure of that transformation, not because every RF pulse is a physically equivalent selector. With dissipation, imperfect preparation, or non-Gaussian dynamics, even that simple determinant need not remain fixed. The transformation class, rather than the general word ``projection,'' determines the invariance claim.\cite{weedbrook,foundation}

\subsection{The stationary source and the ordered laboratory record}
Within \TCGS, an ordered experimental record is a shadow-level registration of a source-complete architecture. That interpretation leaves the observed correlations among pulse duration, phase-space angle, and imaging outcome intact. It also prevents the source from being described as a separate object oscillating at the measured $\omega_k$. The latter frequency belongs to the effective field Hamiltonian and its operational clock calibration.\cite{foundation,logical}

The $2+1$ simulator notation consequently has a valid representational role without deciding the dimensional ontology of Counterspace. The present paper retains the full effective dynamics while refusing the additional identification of a coordinate in that dynamics with a source direction. This is the specific application of the Minkowski-trap distinction relevant to the experiment.\cite{minkowski,foundation}

\section{Quantum probability and the limit of ontological inference}
\label{sec:probability}
\subsection{Readout statistics and source interpretation}
For a prepared effective state $\rho$ and a specified measurement effect $E_A$, the operational probability is
\begin{equation}
 p(A\mid\rho,\mathcal M)=\Tr(\rho E_A).
 \label{eq:born_effect}
\end{equation}
The framework's Born-rule paper interprets such probabilities as projection-invariant shadow assignments rather than source-side chance. The present covariance calculation retains this operational probability calculus. It does not replace it with a classical noise process or derive a new distribution from the BEC images.\cite{born,foundation}

A source-conditioned representation would instead require a space of admissible configurations, a declared measure, and a readout map. Schematically,
\begin{equation}
 p_\sigma(A\mid\mathcal I)
 =\mu\!\left(\{\Xi:\Pi_\sigma(\Xi)\in A\}\mid\mathcal I\right).
 \label{eq:pushforward}
\end{equation}
Here $\Pi_\sigma$ is not the immersion $X_\sigma$, and $\mu$ is not supplied merely by writing an inverse-image symbol. A concrete source model must generate the same operational statistics, including their dependence on state preparation and measurement context.\cite{born,foundation}

The role of Gleason-type reasoning is correspondingly conditional on an event structure with its required additivity and dimensional assumptions. An arbitrary geometric projection does not by itself determine a unique probability measure. In this application the Hilbert-space probabilities and Gaussian covariance are retained as the established domain-level description; the source interpretation is not used to claim that those additional mathematical hypotheses have been derived anew.\cite{born}

\subsection{Uncertainty is not merely an optical disturbance}
The finite ground-state covariance exists before the final imaging pulse in the quantum description. Therefore it would be inaccurate to explain it solely as detector noise or disturbance caused by looking at the atoms. The commutator constrains state preparation as well as readout. Interpreting uncertainty as limited shadow accessibility must preserve that preparation constraint and the state-dependent covariance, not reduce all quantum uncertainty to an imperfect microscope.\cite{zhang2026,weedbrook,born}

This also clarifies the meaning of the visually irregular snapshots. Their reproducible ensemble structure is experimentally real. Their irregularity alone does not tell whether an underlying ontology is fundamentally stochastic, deterministic but nonlocal or contextual, or described in another manner consistent with the same operational probabilities. The present experiment compares quantum-vacuum and thermal covariance models; it does not implement a general discrimination among all interpretations of quantum mechanics.\cite{zhang2026}

\subsection{Gaussian representability and its restricted lesson}
A positive Gaussian covariance can be represented by a classical Gaussian probability density on phase space. For quadratic dynamics, its symmetrically ordered moments propagate by the same linear symplectic transformation. This permits a phase-space calculation of the measured covariance without asserting that the represented phase-space points are literal simultaneous values of all physical observables. Gaussian-state and c-field methods supply the relevant computational setting.\cite{weedbrook,blakie}

The limited lesson is that agreement of these particular second moments cannot identify a unique microscopic outcome ontology. It does not establish a general classical replacement for quantum mechanics, a Bell-local hidden-variable theory, or arbitrary noncontextual value assignments. The retained uncertainty bound and the measurement model remain essential. The \TCGS\ interpretation is disciplined by reproducing those constraints rather than by denying them.\cite{born,foundation}

\section{Identifiability of the carrier and the status of the anchor}
\label{sec:anchor}
\subsection{The observation supports a local distinction}
The relevant observation combines three experimentally specified features: finite transverse confinement, an effective planar spin field, and vacuum-dominated covariance. The first two support the carrier--representation distinction through Equations~\eqref{eq:width} and \eqref{eq:g2d}. The third identifies the quantum state regime in which the planar field was imaged. The vacuum spectrum is not needed to calculate the oscillator width, and the finite-width calculation is not needed to derive the quantum $1/2$ contribution. Their conjunction is informative precisely because the two arguments remain independent.\cite{zhang2026}

\begin{proposition}[Planar-data non-identifiability in a reduced model]
Within the common-profile, weak-contact reduction, consider two atomic realizations with the same nominal mass, areal density, coherent coupling, and planar geometry. If all scattering lengths and the transverse oscillator length are transformed as
\begin{equation}
 a_{\varsigma\varsigma'}\longmapsto\lambda a_{\varsigma\varsigma'},
 \qquad \lop\longmapsto\lambda\lop,
 \qquad \lambda>0,
 \label{eq:degeneracy}
\end{equation}
then the effective interaction coefficients, reduced spin dispersion, and equilibrium spin covariance are unchanged, although the transverse widths differ.
\end{proposition}
\begin{proof}
Equation~\eqref{eq:g2d} depends on each ratio $a_{\varsigma\varsigma'}/\lop$, which is unchanged. Consequently $\kappa$, $g_d$, $Z_0$, and $\musp$ are unchanged at fixed areal density. Equations~\eqref{eq:ab}--\eqref{eq:field_variances} then give the same reduced predictions, while Equation~\eqref{eq:width} gives a width multiplied by $\lambda$.
\end{proof}

The proposition is a non-identifiability construction within a model class, not a claim that all three scattering lengths can be tuned together arbitrarily in the reported potassium apparatus. Both realizations must remain inside the regime where the reduction is valid. The construction shows why complete knowledge of the planar coefficients need not reconstruct the transverse carrier without additional information.

\subsection{Recovering the suppressed coordinate through a form factor}
For a probe sensitive to transverse momentum transfer $q_z$, the normalized density profile has form factor
\begin{equation}
 F_z(q_z)=\int_{\R}|f_0(z)|^2e^{\ii q_zz}\,\dd z
 =\exp\!\left(-\frac{q_z^2\lop^2}{4}\right),
 \qquad |F_z(q_z)|^2=e^{-q_z^2\lop^2/2}.
 \label{eq:formfactor}
\end{equation}
This is a prediction for the nominal Gaussian profile and a specified coherent profile-sensitive channel. A real scattering experiment must additionally include internal transitions, probe geometry, and its detection response. The formula is not attributed to a transverse measurement performed in the vacuum-imaging paper.

At fixed $q_z$, the two realizations in Equation~\eqref{eq:degeneracy} generally have different form factors. A sufficiently calibrated transverse-sensitive measurement can therefore recover information absent from the planar spectrum. The inverse problem is improved by adding an observable sensitive to the discarded variable, not by interpreting a planar fit more strongly than its information content allows.

\subsection{Empirical anchoring and comparative confirmation}
Let two interpretations assign the same likelihood to the available data $D$ under the same controlled apparatus model:
\begin{equation}
 p(D\mid\mathcal H_1)=p(D\mid\mathcal H_2).
 \label{eq:equal_likelihood}
\end{equation}
Then their likelihood ratio is unity for those data. The observation can be an excellent illustration of a distinction common to both interpretations without selecting one interpretation over the other. This is a logical consequence of equal predictions, not a judgement against either ontology.

Accordingly, the BEC result is an additional empirical anchor for the \TCGS\ carrier--representation distinction, while its immediate quantitative explanation is also supplied by ordinary confined-atom theory. A future source-derived relation could discriminate more strongly only if it predicted an independently specified observable not already fixed by the retained atomic Hamiltonian. The foundational programme and the conditional Planck analysis both require that distinction between governing architecture and a testable domain-level map.\cite{foundation,planck}

\section{Comparison with spectral reduction and holographic descriptions}
\label{sec:comparisons}
\subsection{Dimensional reduction is not one universal operation}
In a confined gas, reduction proceeds by suppressing or eliminating transverse excitations. In a diffusion-based dimensional estimator, reduction concerns a scale-dependent spectral law. In a bulk--boundary duality, it concerns a correspondence between theories or observable structures. These operations can share a lower effective dimensional label while preserving different information and requiring different hypotheses. The dimensionality literature distinguishes the estimators precisely for this reason.\cite{carlip}

For a positive diffusion generator $\mathcal D$, define a return quantity $\mathcal P(s_D)$ and its spectral dimension by
\begin{equation}
 d_s(s_D)=-2\frac{\dd\log\mathcal P(s_D)}{\dd\log s_D}.
 \label{eq:spectral_dimension}
\end{equation}
The diffusion parameter $s_D$ is not automatically a laboratory clock or a source coordinate. The numerical value of $d_s$ depends on the operator, measure, ensemble, and scale. It must not be substituted for the topological dimension of a carrier without a further theorem.\cite{carlip,planck}

Causal dynamical triangulations provide a numerical example of scale-dependent spectral dimension, while asymptotically safe gravity provides another construction with ultraviolet dimensional reduction. These are technical results in specified quantum-gravity models, not measurements that a cold atomic sample loses a spatial direction. Their relevance here is to illustrate the necessity of naming the dimensional estimator rather than rejecting dimensional reduction in general.\cite{cdt,reuter}

Ho\v rava's anisotropic continuum calculation makes the distinction explicit. For $D$ spatial coordinates and scaling exponent $z$, the result $d_s=1+D/z$ gives $d_s=2$ when $D=z=3$, without replacing the underlying smooth coordinate manifold by a two-dimensional one. The cold-atom and spectral examples therefore share a non-identification principle, not a common microscopic mechanism.\cite{horava}

\subsection{A boundary can be two-dimensional while its domain remains three-dimensional}
The \TCGS\ Planck analysis considers a regular source level set $\Gamma=B_C^{-1}(1)$ and a rank-three immersion $X:\Sigma^3\hookrightarrow C^4$ transverse to $\Gamma$. Its preimage $X^{-1}(\Gamma)$, when nonempty, is two-dimensional, while the associated admissible domain in $\Sigma$ remains three-dimensional. The theorem is a differential-topological statement, with scale and physical interpretation requiring separate bridges.\cite{planck}

The BEC result is not an instance of that preimage theorem. Harmonic confinement, a mode projection, and a regular level-set preimage are different constructions. Nevertheless, both analyses block the same invalid inference: the appearance of a lower-dimensional mathematical object does not by itself determine the dimensionality of the physical carrier or its source. The restriction matters because an analogy should preserve the kind of statement established, not only the numbers appearing in it.\cite{planck}

\subsection{Holography must be evaluated in its own category}
The holographic principle concerns entropy bounds and the organization of degrees of freedom; the original literature is not simply a proposal that a laboratory material becomes a zero-thickness sheet. AdS/CFT supplies a particular duality between specified theories, with a dictionary relating their observables. These technical claims are categorically different from the single-mode reduction of a harmonically confined gas.\cite{susskind,maldacena,bousso}

Therefore, finite atomic thickness is not an experimental refutation of holographic duality. Conversely, the existence of a useful boundary description would not establish that the atomic carrier in this experiment is literally two-dimensional. The proper question is what each representation retains and whether the necessary source-to-observable bridge is supplied. A conclusion about one class cannot be transferred to another merely because both use lower-dimensional variables.

The \TCGS\ alternative is stated positively through 4D-Counterduction: a source-grounded constitutive realization of a physical three-dimensional shadow. It assumes neither a literal holographic boundary nor a globally invertible encoding. The present empirical anchor belongs to that source--carrier--representation hierarchy without making the atomic planar field a miniature physical model of the full Counterspace relation.\cite{foundation}

\section{Invariants, artifacts, assumptions, and interpretations}
\label{sec:ledger}
The analysis distinguishes what is observed, what is derived under explicit hypotheses, and what is interpreted through the framework. The word invariant is always relative to a declared comparison class. A mode covariance can have a symplectic invariant; a calibrated dimensionless spectral curve can transfer between matched experiments; a source-level invariant requires the appropriate selector comparison. These are not interchangeable achievements.\cite{foundation}

\begin{longtable}{@{}L{0.25\textwidth}L{0.27\textwidth}L{0.40\textwidth}@{}}
\caption{Claim-status ledger for the present application.\label{tab:ledger}}\\
\toprule Claim or object & Status & Scope or requirement\\\midrule
\endfirsthead
\toprule Claim or object & Status & Scope or requirement\\\midrule
\endhead
\bottomrule\endfoot
Vacuum-dominated spin spectrum & Reported observation and model comparison & Source temperature fit, mode selection, and detection correction retained.\cite{zhang2026}\\
Finite transverse confinement & Reported apparatus property & Establishes a physical confinement scale, not a measured universal minimum length.\cite{zhang2026}\\
$\lop\simeq0.509\,\mu\mathrm m$ & Derived nominal value & Harmonic profile, approximate isotope mass, and reported trap frequency.\\
$g_{2D}\propto\lop^{-1}$ & Established reduction & Common transverse mode and weak-contact approximation.\cite{hadzibabic}\\
Zero-width finite-resource obstruction & Derived model-level result & Canonical transverse pair, positive atomic mass, and finite kinetic-energy bound.\\
$\det C=1/4$ for a ground-state mode & Canonical prediction & Not reconstructed from amplitude difference alone.\\
$A_Y\propto\omega^{-1}$ & Vacuum-limit prediction & Per-mode covariance scaling; not a spatial-dimensional estimator.\\
$\eta\simeq0.7$ & Reported calibration factor & Imperfect detection is the authors' likely explanation, not uniquely identified here.\cite{zhang2026}\\
Filtered image morphology & Processing-dependent record & Must retain window, masking, transfer function, and ensemble definition.\\
Three-dimensional shadow carrier & Governing framework architecture & Not refitted from each cold-atom spectrum.\cite{foundation}\\
Source-side deterministic completeness & Framework interpretation & Must reproduce the operational probability and covariance; not selected by these second moments alone.\cite{born}\\
Full source-to-spin map & Unexecuted domain bridge & Requires explicit construction beyond relabeling the standard atomic Hamiltonian.\\
Transverse form-factor test & Prospective model-level test & Requires a calibrated transverse-sensitive probe; not already performed in the source.\\
Refutation of holographic duality & Unsupported by this experiment & A confined atomic field and a bulk--boundary duality are different categories.\\
\end{longtable}

A practical consequence of the ledger is that physical protocol dependence is not automatically a defect. Changing the state preparation should change the covariance. Changing the detector response should change the raw image. The scientific task is to identify the appropriate corrected relation and its declared stability, not to demand that every measurable quantity be invariant under every intervention.

Equally, framework interpretation is not a substitute for constructing the maps. A label such as source-conditioned covariance does not specify the kernel, normalization, state preparation, or probability measure. The existing oscillator calculation provides the operational target that any more fundamental map must recover, including its finite-temperature and finite-readout corrections.\cite{foundation,zhang2026}

\section{A transferable experimental and analytical programme}
\label{sec:tests}
\subsection{Confinement variation with independently measured inputs}
At fixed mass and three-dimensional scattering length, the nominal reduction predicts
\begin{equation}
 \lop\propto\omega_z^{-1/2},\qquad
 \sigma_z\propto\omega_z^{-1/2},\qquad
 g_{2D}\propto\omega_z^{1/2}.
 \label{eq:confinement_prediction}
\end{equation}
A controlled sweep of transverse frequency should therefore be accompanied by independent calibration of trap shape, scattering lengths, areal density, RF parameters, and imaging response. Otherwise a width-induced interaction change could be hidden inside a refitted $\mu_s$ and the carrier dependence would not have been tested.

A useful test compares the measured interaction change with the width inferred from an independent axial measurement, rather than fitting both from the same planar spectrum. Agreement would support the shared-profile reduction; systematic disagreement would motivate excited-mode, non-Gaussian-profile, or scattering corrections. It would not automatically imply a failure of the physical carrier or the appearance of a new source sector.

\subsection{Dimensionless collapse across matched planar sectors}
Define
\begin{equation}
 u_k=\frac{\epsilon_k}{\musp},\qquad
 v=\frac{\hbar\Omega'}{\musp},\qquad
 w_k=\frac{\hbar\omega_k}{\musp}
     =\sqrt{(u_k+v)(u_k+v+2)}.
 \label{eq:dimensionless}
\end{equation}
In the ideal vacuum-dominated limit,
\begin{equation}
 w_k A_Y(k)=1.
 \label{eq:vacuum_collapse}
\end{equation}
At finite temperature it becomes $w_kA_Y=2N_k+1$. With a known multiplicative detection attenuation, the corrected amplitude must be used; a separately fitted factor for every target condition would not constitute a transferred prediction.

This test is motivated by the source's use of a common fluctuation law across wavevectors and coherent couplings. Extending it across deliberately different transverse confinement conditions would probe whether the same reduced coefficients account for the carrier change. It tests an effective-theory matching relation rather than a new universal spectral law.\cite{zhang2026}

\subsection{Full covariance tomography}
At a minimum, three linearly independent quadrature orientations are required to reconstruct a real symmetric $2\times2$ covariance once absolute offsets are calibrated. More angles allow residual and harmonic-content checks. The analysis should recover the cross covariance rather than fixing it to zero solely because the preparation is intended to be stationary.

The resulting checks are
\begin{equation}
 C\succeq0,\qquad \det C\geq\frac14,
 \qquad C_{QP}\simeq0\;\text{for a calibrated equilibrium mode}.
 \label{eq:tomography_checks}
\end{equation}
A result apparently below the uncertainty limit should first trigger a units, normalization, loss, or subtraction audit. A result above it could indicate thermal occupation, mixing, excess technical noise, or incomplete mode resolution. Only a complete forward model can separate these possibilities.\cite{weedbrook}

For multimode tomography, mode mixing must also be included. A one-mode covariance determinant is not invariant under tracing over correlations with other modes. A projected imaging mode can therefore appear mixed even when the larger many-body state is pure. This is another reason not to infer source randomness or source incompleteness directly from one reduced covariance.\cite{weedbrook,foundation}

\subsection{Axial response as a carrier discriminator}
A transverse-sensitive probe should test Equation~\eqref{eq:formfactor} or a more complete experimentally derived response. The relevant discriminator is dependence on $q_z\lop$, not merely a new image in the same column-integrated geometry. In particular, selecting two realizations with matched planar coefficients but different independently measured widths would directly test the non-identifiability construction of Section~\ref{sec:anchor}.

Excited transverse modes provide another discriminator. As the relevant energy, temperature, or coupling matrix elements increase, a common single-mode profile may cease to suffice. The correct reduced theory then acquires additional mode populations or effective corrections. Tracking this crossover tests the boundary of the approximation while leaving the three-dimensional carrier explicit throughout.

\subsection{Held-out calibration and thermal discrimination}
The detection efficiency and the temperature affect the amplitude differently. In a vacuum-dominated range, an unknown multiplicative efficiency can be almost degenerate with the overall normalization. Thermal occupation introduces frequency-dependent curvature through Equation~\eqref{eq:coth}. A broad, independently selected frequency range is therefore more informative than a narrow range over which both effects resemble a scale change.

A defensible protocol fixes the admissible frequency window, ramp criterion, imaging correction, and uncertainty model before fitting the temperature. Calibration runs determine shared nuisance parameters; held-out runs assess their transfer. Residuals should be reported against wavevector, coupling, density, transverse frequency, and readout angle. A model that requires an arbitrary efficiency function of all those variables has not isolated a universal fluctuation law.

\subsection{What would make a source-derived contribution distinct}
The tests above are deliberately rooted in known confined-atom and Gaussian-state physics. Their purpose is to operationalize the carrier--representation distinction. A specifically source-derived extension would need an independently constructed relation between admissible source data and an observable not already fixed by those retained equations. Its parameters would have to be set outside the target observations, and its additional prediction would need a stated failure condition.\cite{foundation,planck}

For example, a proposed source-induced correction to a cross-mode covariance would need a defined tensor or kernel, a dimensional normalization, an experimental estimator, and a comparison against finite-temperature, imaging, and interaction corrections. Merely identifying a residual with Counterspace would not supply any of these. The framework becomes more restrictive, not less, when the standard experimental description is retained in full.

\section{Discussion}
\label{sec:discussion}
\subsection{What is gained by the carrier-first interpretation}
The principal gain is an explicit separation of explanatory levels. The experiment prepares a physical atomic carrier, selects an effective field sector, applies a readout operation, and reconstructs covariance statistics. Each step has its own mathematical objects and possible errors. Treating the final field image as a literal picture of an autonomous vacuum would skip the apparatus and reduction maps that make the observation possible.\cite{zhang2026}

A carrier-first account does not diminish the quantum achievement. The vacuum-dominated scaling remains a nontrivial result, especially because the field modes are resolved spatially and compared across scales. The account instead identifies which inference is specifically dimensional: the transverse coordinate remains in the carrier profile and effective coupling even when the recorded field depends only on two coordinates.\cite{zhang2026,hadzibabic}

The resulting anchor is particularly suited to the \TCGS\ map--territory distinction because the suppressed variable is not merely philosophical. It has a finite width, a localization cost, and a possible form-factor signature. The reduced map is empirically successful precisely while remaining non-exhaustive of the carrier. Within the framework, that is the relevant downstream example of three-dimensional shadow realization.\cite{foundation}

\subsection{What is not novel and what the present synthesis contributes}
Harmonic confinement, quasi-two-dimensional coupling, oscillator uncertainty, Gaussian covariance propagation, and the distinction between thermal and vacuum occupation are established physics. They are not claimed as discoveries of \TCGS. The source experiment itself supplies the vacuum-spectrum formula and the apparatus model. The present contribution is their integration into a typed carrier analysis, together with an explicit local irreducibility statement, a covariance-identifiability audit, and a set of complementary carrier tests.\cite{zhang2026,hadzibabic,petrov,weedbrook}

The strongest conceptual result is not simply that ``the atoms are 3D.'' It is that reduced observables can be exactly represented on a lower-dimensional effective sector while the physical realization remains nontrivially dependent on omitted structure. Equations~\eqref{eq:expectation_lift} and \eqref{eq:degeneracy} exhibit this distinction mathematically. Equation~\eqref{eq:formfactor} identifies how a different measurement can recover information absent from the planar sector.

The covariance analysis adds a parallel lesson. Anisotropy is observable, a determinant is invariant under a specified class of canonical transformations, and a thermal occupation is inferred only after a state model is supplied. These are different levels of information. The analogy with carrier reduction is useful because both cases show how a successful observable need not be a complete inverse of its realization.

\subsection{Limits of the present analysis}
The nominal harmonic width has not been independently measured here. The source's rounded trap frequency and isotope label determine a model value, while actual axial populations, anharmonicity, interactions, and state-dependent confinement could modify the profile. The nonrelativistic finite-energy bound remains valid under its stated canonical assumptions but does not constitute a universal theorem about every proposed fundamental ontology.

The statistical assessment is likewise limited to the published equations, methods, figures, and reported fit. No raw-image likelihood, independent posterior for $T$, or revised detection efficiency is claimed. The source's residual efficiency factor is retained as a calibration issue, and the proposed tomography and axial-response experiments remain prospective.\cite{zhang2026}

At framework level, no explicit global source configuration, constitutive kernel, or source-to-spin Hamiltonian has been computed. Counterspace and 4D-Counterduction organize the interpretation, but their use does not turn the standard oscillator calculation into an independent reconstruction of the complete source. The source-specific closure burden is the derivation of a transferable map, not the repetition of the ontology in place of a calculation.\cite{foundation,planck}

Finally, neither the finite-width result nor the covariance analysis overturns the technical use of lower-dimensional field theories, spectral dimensional reduction, or holographic duality. Each remains subject to its own definitions, derivations, and empirical scope. The target of criticism is the promotion of a successful representation into an unearned claim of ontological identity.\cite{carlip,bousso,foundation}

\section{Conclusion}
\label{sec:conclusion}
Vacuum-fluctuation imaging in a planar two-component BEC provides a concrete empirical anchor for the distinction between effective field dimension and physical carrier dimension. The apparatus has finite transverse confinement, its nominal harmonic state has $\lop\simeq0.509\,\mu\mathrm m$ and $\sigma_z\simeq0.360\,\mu\mathrm m$, and its planar interaction retains the suppressed direction through $g_{2D}\propto\lop^{-1}$. Within the massive-atom quantum model, exactly vanishing transverse width cannot be reached at uniformly bounded kinetic energy.\cite{zhang2026}

Independently, the spin covariance reproduces the vacuum-dominated inverse-frequency amplitude. The measured quadrature anisotropy is not a full covariance determinant, the per-mode exponent is not a spatial-dimensional estimator, and the sequence of destructive measurements is not a movie of an ontic vacuum fluctuating in source time. These distinctions preserve the result's quantum content while delimiting its interpretation.\cite{zhang2026}

Under \TCGS, the appropriate conclusion is therefore specific: the experiment realizes a lower-dimensional effective description without converting its three-dimensional shadow carrier into an autonomous lower-dimensional substrate. It exemplifies a framework distinction whose governing source architecture remains upstream of this application. Transverse-sensitive measurements, matched-confinement comparisons, and full covariance tomography can sharpen the empirical bridge without confusing it with a universal source reconstruction.\cite{foundation,logical}

The general lesson is not that reduced maps are inadequate because they are reduced. It is that their success should be interpreted at the level of the information they preserve. Here the missing spatial direction is not abolished; it remains in the physical conditions that make the two-dimensional field readable.

\appendix
\section{Finite-initial-coupling quench calculation}
\label{app:quench}
Let an initial centred thermal covariance be
\begin{equation}
 C_i=s_i\diag(r_i,r_i^{-1}),\qquad s_i=N_i+\frac12.
 \label{eq:Ci}
\end{equation}
After a sudden change to the final Hamiltonian, the propagator is
\begin{equation}
 S_f(t)=\begin{pmatrix}
 c&r_fs\\-s/r_f&c
 \end{pmatrix},\qquad
 c=\cos\omega_ft,\quad s=\sin\omega_ft.
 \label{eq:Sf}
\end{equation}
The covariance is $C(t)=S_f(t)C_iS_f(t)^T$, so
\begin{align}
 C_{QQ}(t)&=s_i\left(r_ic^2+\frac{r_f^2}{r_i}s^2\right),\nonumber\\
 C_{PP}(t)&=s_i\left(\frac{r_i}{r_f^2}s^2+\frac{c^2}{r_i}\right),\nonumber\\
 C_{QP}(t)&=s_i\left(\frac{r_f}{r_i}-\frac{r_i}{r_f}\right)cs.
 \label{eq:quench_covariance}
\end{align}
Direct multiplication gives $\det C(t)=s_i^2$, as required by symplectic propagation. The changing diagonal variances therefore reflect redistribution of covariance, not an increase of the covariance determinant under this ideal unitary quench.

The peak-to-peak $Q$ variance is
\begin{equation}
 A_Q=s_i\left|\frac{r_f^2}{r_i}-r_i\right|.
 \label{eq:quench_amplitude}
\end{equation}
For a deep-Rabi initial vacuum, $r_i\to1$ and $s_i=1/2$. Since $r_f^2-1=2\musp/a_f$, the result becomes
\begin{equation}
 A_Q=\frac{\musp}{a_f}
 =\frac{\musp}{\epsilon_k+\hbar\Omega'_f}.
 \label{eq:quench_source}
\end{equation}
This is the ideal structure-factor amplitude used for the source's amplification protocol. A finite $r_i$ changes the initial covariance before any readout correction is considered. The experiment's quoted finite-parameter correction additionally includes the actual pulse operation and is not inferred from Equation~\eqref{eq:quench_amplitude} alone.\cite{zhang2026}

For matched initial and final Hamiltonians, $r_i=r_f$, Equation~\eqref{eq:quench_amplitude} vanishes. The ground-state covariance is then stationary. This limiting case is an internal check that nonzero quantum variance does not itself require a breathing structure factor.

\section{Projected dynamics and an explicit error hypothesis}
\label{app:projection}
The mode expansion defines a projector $P$ onto a chosen transverse sector and $Q=I-P$ onto its complement. For an eigenvalue equation $H\psi=E\psi$, suppose that $E-QHQ$ is invertible on the relevant complementary domain. Eliminating $Q\psi$ gives
\begin{equation}
 H_{\mathrm{eff}}(E)
 =PHP+PHQ(E-QHQ)^{-1}QHP.
 \label{eq:schur}
\end{equation}
This follows by solving $(E-QHQ)Q\psi=QHP\psi$ and substituting into the projected equation. The second term is the virtual-sector correction; setting it to zero is an approximation, not a geometric deletion of the complementary space.

For a quantitative estimate, restrict to a specified energy sector on which the coupling blocks are bounded. If
\begin{equation}
 \|PHQ\|\leq v,\qquad
 \|(E-QHQ)^{-1}\|\leq\Delta^{-1},\qquad \Delta>0,
 \label{eq:projection_assumptions}
\end{equation}
then
\begin{equation}
 \|H_{\mathrm{eff}}(E)-PHP\|\leq\frac{v^2}{\Delta}.
 \label{eq:projection_error}
\end{equation}
The estimate is conditional on a real spectral separation and bounded restricted couplings. An infinite-dimensional many-body Hamiltonian does not inherit a global bound of this form merely because the one-particle harmonic oscillator has a gap.

In the BEC setting, density-dependent interactions, internal-state operations, and possible excited transverse populations determine whether such a restricted-sector estimate is useful. The published nominal gap and density scale motivate the reduction, but the complete $v$ and $\Delta$ needed for a rigorous many-body bound are not supplied here. The purpose of the construction is to make the missing quantitative condition explicit rather than replace it with an unrestricted assertion of exact two-dimensionality.

\section{Confinement families, uncertainty propagation, and numerical checks}
\label{app:numeric}
For a harmonic frequency family $\omega_z(q)=q\omega_{z,0}$ with $q>0$,
\begin{equation}
 \lop(q)=q^{-1/2}\lop(1),\quad
 \sigma_z(q)=q^{-1/2}\sigma_z(1),\quad
 K_z(q)=qK_z(1).
 \label{eq:q_family}
\end{equation}
At fixed scattering length, $g_{2D}(q)=q^{1/2}g_{2D}(1)$. This family is an analytic sensitivity test. It is not an experimental sweep performed for the present paper.

\begin{table}[htbp]
\centering\small
\caption{Illustrative harmonic confinement family at fixed nominal $^{39}\mathrm K$ mass. All rows are calculated, not new measurements.}
\label{tab:family}
\begin{tabular}{@{}rrrrr@{}}
\toprule
$\omega_z/(2\pi)$ (kHz)&$\lop$ ($\mu$m)&$\sigma_z$ ($\mu$m)&$K_z/k_B$ (nK)&$g_{2D}/g_{2D,1\mathrm{kHz}}$\\\midrule
0.5&0.720&0.509&6.00&0.707\\
1.0&0.509&0.360&12.0&1.000\\
2.0&0.360&0.255&24.0&1.414\\
4.0&0.255&0.180&48.0&2.000\\
\bottomrule
\end{tabular}
\end{table}

For jointly uncertain $m$ and $\omega_z$, the first-order logarithmic variance is
\begin{equation}
 \Var(\log\lop)=\frac14\left[
 \Var(\log m)+\Var(\log\omega_z)
 +2\Cov(\log m,\log\omega_z)\right].
 \label{eq:width_uncertainty}
\end{equation}
At fixed mass, the corresponding interaction uncertainty is
\begin{equation}
 \Var(\log g_{2D})
 =\Var(\log a)+\frac14\Var(\log\omega_z)
 +\Cov(\log a,\log\omega_z).
 \label{eq:coupling_uncertainty}
\end{equation}
The covariance terms are not optional when the same control field enters multiple calibrations. The equations also show why a fitted interaction scale should not be reused as an independent determination of the width without accounting for shared information.

The companion numerical script evaluates the harmonic scales, the finite-energy saturation identity, the dimensionless dispersion, the covariance determinant, and the quench-amplitude limit. It generates only derived values and analytic test cases. No generated row is labeled as an experimental image, an observed spectrum, or a newly fitted temperature.

\section{A minimal source-to-observable closure checklist}
\label{app:closure}
A concrete \TCGS\ completion of this atomic sector would require more than the source ontology and a successful oscillator analogy. First, the source configuration and its admissible selector family must be specified. Second, the map to an operational spin state must be constructed with its normalization. Third, the effective Hamiltonian and its transverse dependence must be recovered. Fourth, the registered measurement probabilities must agree with the same preparation and readout model used experimentally. Finally, any additional source-derived relation must transfer to held-out conditions without independent retuning.\cite{foundation}

The relevant diagram is a composition of distinct maps:
\begin{equation}
 \begin{aligned}
 \Xi&\xrightarrow{\mathrm{Ctd}^{\sigma}_{4,S,G}}
 \text{physical atomic readout},\\
 \text{physical atomic readout}
 &\xrightarrow{\text{effective reduction}}\rho_{\mathrm{spin}}
 \xrightarrow{\text{measurement}}p(\mathrm{OD},k,\vartheta).
 \end{aligned}
 \label{eq:closure_chain}
\end{equation}
The displayed atomic and statistical arrows are not already a construction of the first arrow. Conversely, any claimed first-arrow construction must reproduce the empirical consequences of the later ones. This is the framework's source--map accountability applied to the vacuum-fluctuation experiment.\cite{foundation,planck}

\section*{Data, code, and scope statement}
No new experimental data are reported. Source-reported apparatus values and fitted quantities are cited to their original publications. The accompanying numerical analysis evaluates equations and nominal-scale calculations in this article. The manuscript is an analytical application of \TCGS\ and a critical synthesis of the cited atomic-field and dimensionality literature; it is not an independent experimental replication.

\section*{Declaration of perspective}
The author develops the \TCGS\ framework. Framework-governed interpretation is distinguished throughout from source-reported observations and from the conventional quantum-mechanical calculations used to evaluate the laboratory carrier.

\section*{License}
Copyright \textcopyright\ 2026 Henry Arellano-Pe\~na. This manuscript is distributed under the Creative Commons Attribution 4.0 International License (CC BY 4.0), \url{https://creativecommons.org/licenses/by/4.0/}. The license applies to this manuscript and its original analytical material, not automatically to third-party publications cited or supplied with it.

\begin{thebibliography}{99}
\small
\setlength{\itemsep}{0pt}
\setlength{\parskip}{0pt}
\bibitem{zhang2026} Y. Zhang, F. Wang, Y. Jiang, A. C. Jenkins, P. H. C. Wong, C. Eigen, G. Carlse, and Z. Hadzibabic, ``Imaging the vacuum fluctuations of a quantum field,'' \arxiv{2608.20311v1} (20 August 2026).
\bibitem{foundation} H. Arellano-Pe\~na, ``Timeless Counterspace \& Shadow Gravity---A Unified Framework: 4D-Counterduction, Foundational Consistency, and Cartographic Inquiry,'' \emph{Preprints} (11 August 2026), version 2. \doi{10.20944/preprints202511.1472.v2}.
\bibitem{logical} H. Arellano-Pe\~na, ``Why 4D-Counterspace Exists: A Formal Logical Demonstration within the TCGS--SEQUENTION Framework,'' preprint (28 April 2026). \doi{10.13140/RG.2.2.29353.76641}.
\bibitem{planck} H. Arellano-Pe\~na, ``Planck-Scale Dimensional Reduction Without Topological Collapse: Spectral Dimension, Curvature Thresholds, and a Conditional TCGS--SEQUENTION Interpretation,'' preprint (1 September 2026).
\bibitem{hadzibabic} Z. Hadzibabic and J. Dalibard, ``Two-dimensional Bose fluids: An atomic physics perspective,'' \emph{Rivista del Nuovo Cimento} \textbf{34}, 389 (2011). \doi{10.1393/ncr/i2011-10066-3}; \arxiv{0912.1490}.
\bibitem{petrov} D. S. Petrov, M. Holzmann, and G. V. Shlyapnikov, ``Bose-Einstein Condensation in Quasi-2D Trapped Gases,'' \emph{Physical Review Letters} \textbf{84}, 2551--2555 (2000). \doi{10.1103/PhysRevLett.84.2551}; \arxiv{cond-mat/9909344}.
\bibitem{chomaz} L. Chomaz, L. Corman, T. Bienaim\'e, R. Desbuquois, C. Weitenberg, S. Nascimb\`ene, J. Beugnon, and J. Dalibard, ``Emergence of coherence via transverse condensation in a uniform quasi-two-dimensional Bose gas,'' \emph{Nature Communications} \textbf{6}, 6162 (2015). \doi{10.1038/ncomms7162}; \arxiv{1411.3577}.
\bibitem{navon} N. Navon, R. P. Smith, and Z. Hadzibabic, ``Quantum gases in optical boxes,'' \emph{Nature Physics} \textbf{17}, 1334--1341 (2021). \doi{10.1038/s41567-021-01403-z}; \arxiv{2106.09716}.
\bibitem{precursor} Y. Zhang, F. Wang, P. H. C. Wong, A. C. Jenkins, Y. Jiang, K. Konstantinou, G. Carlse, N. Dogra, J. H. Thywissen, C. Eigen, and Z. Hadzibabic, ``Analog Simulation of Massive Relativistic Quantum Fields in $2+1$ Dimensions,'' \arxiv{2603.08840v2} (20 August 2026).
\bibitem{son} D. T. Son and M. A. Stephanov, ``Domain walls of relative phase in two-component Bose-Einstein condensates,'' \emph{Physical Review A} \textbf{65}, 063621 (2002). \doi{10.1103/PhysRevA.65.063621}; \arxiv{cond-mat/0103451}.
\bibitem{hung} C.-L. Hung, X. Zhang, L.-C. Ha, S.-K. Tung, N. Gemelke, and C. Chin, ``Extracting density-density correlations from in situ images of atomic quantum gases,'' \emph{New Journal of Physics} \textbf{13}, 075019 (2011). \doi{10.1088/1367-2630/13/7/075019}; \arxiv{1105.0030}.
\bibitem{blakie} P. B. Blakie, A. S. Bradley, M. J. Davis, R. J. Ballagh, and C. W. Gardiner, ``Dynamics and statistical mechanics of ultra-cold Bose gases using c-field techniques,'' \emph{Advances in Physics} \textbf{57}, 363--455 (2008). \doi{10.1080/00018730802564254}; \arxiv{0809.1487}.
\bibitem{weedbrook} C. Weedbrook, S. Pirandola, R. Garc\'ia-Patr\'on, N. J. Cerf, T. C. Ralph, J. H. Shapiro, and S. Lloyd, ``Gaussian quantum information,'' \emph{Reviews of Modern Physics} \textbf{84}, 621--669 (2012). \doi{10.1103/RevModPhys.84.621}; \arxiv{1110.3234}.
\bibitem{benea} I.-C. Benea-Chelmus, F. F. Settembrini, G. Scalari, and J. Faist, ``Electric field correlation measurements on the electromagnetic vacuum state,'' \emph{Nature} \textbf{568}, 202--206 (2019). \doi{10.1038/s41586-019-1083-9}; \arxiv{1809.01785}.
\bibitem{moskalenko} A. S. Moskalenko, C. Riek, D. V. Seletskiy, G. Burkard, and A. Leitenstorfer, ``Paraxial Theory of Direct Electro-optic Sampling of the Quantum Vacuum,'' \emph{Physical Review Letters} \textbf{115}, 263601 (2015). \doi{10.1103/PhysRevLett.115.263601}; \arxiv{1508.06953}.
\bibitem{minkowski} H. Arellano-Pe\~na, ``Escaping the Minkowski Trap: Why Time Cannot Be a Dimension,'' preprint, manuscript dated 9 June 2026.
\bibitem{born} H. Arellano-Pe\~na, ``From Deterministic Counterspace to Stochastic Shadow: Deriving the Born Rule as a Projection Invariant in TCGS-SEQUENTION,'' \emph{Preprints} (26 December 2025), version 1. \doi{10.20944/preprints202512.2392.v1}.
\bibitem{analogue} C. Barcel\'o, S. Liberati, and M. Visser, ``Analogue Gravity,'' \emph{Living Reviews in Relativity} \textbf{14}, 3 (2011). \doi{10.12942/lrr-2011-3}; \arxiv{gr-qc/0505065v3}.
\bibitem{carlip} S. Carlip, ``Dimension and Dimensional Reduction in Quantum Gravity,'' \emph{Classical and Quantum Gravity} \textbf{34}, 193001 (2017). \doi{10.1088/1361-6382/aa8535}; \arxiv{1705.05417}.
\bibitem{cdt} J. Ambj\o rn, J. Jurkiewicz, and R. Loll, ``The Spectral Dimension of the Universe is Scale Dependent,'' \emph{Physical Review Letters} \textbf{95}, 171301 (2005). \doi{10.1103/PhysRevLett.95.171301}; \arxiv{hep-th/0505113}.
\bibitem{horava} P. Ho\v rava, ``Spectral Dimension of the Universe in Quantum Gravity at a Lifshitz Point,'' \emph{Physical Review Letters} \textbf{102}, 161301 (2009). \doi{10.1103/PhysRevLett.102.161301}; \arxiv{0902.3657}.
\bibitem{reuter} O. Lauscher and M. Reuter, ``Fractal Spacetime Structure in Asymptotically Safe Gravity,'' \emph{Journal of High Energy Physics} \textbf{2005}(10), 050 (2005). \doi{10.1088/1126-6708/2005/10/050}; \arxiv{hep-th/0508202}.
\bibitem{susskind} L. Susskind, ``The World as a Hologram,'' \emph{Journal of Mathematical Physics} \textbf{36}, 6377--6396 (1995). \doi{10.1063/1.531249}; \arxiv{hep-th/9409089}.
\bibitem{maldacena} J. M. Maldacena, ``The Large N Limit of Superconformal Field Theories and Supergravity,'' \emph{Advances in Theoretical and Mathematical Physics} \textbf{2}, 231--252 (1998). \doi{10.4310/ATMP.1998.v2.n2.a1}; \arxiv{hep-th/9711200}.
\bibitem{bousso} R. Bousso, ``The Holographic Principle,'' \emph{Reviews of Modern Physics} \textbf{74}, 825--874 (2002). \doi{10.1103/RevModPhys.74.825}; \arxiv{hep-th/0203101}.
\end{thebibliography}
\end{document}
